%This chapter was modified on 6-4-05.
%\newcommand{\NA}{{\rm NA}}
%\setcounter{chapter}{8}
\chapter {Teorema Limit Pusat}\label{chp 9} 
\section[Percobaan Bernoulli]{Teorema Limit Pusat untuk Percobaan Bernoulli}\label{sec 9.1}
Teorema dasar kedua dalam probabilitas adalah \emx {Teorema Limit
Pusat.}\index{Teorema Limit Pusat}  Teorema ini menyatakan bahwa jika $S_n$ merupakan jumlah $n$
peubah acak yang saling bebas, maka fungsi distribusi $S_n$ dapat
diaproksimasi dengan baik oleh jenis fungsi kontinu tertentu yang dikenal sebagai fungsi densitas
normal, yang diberikan oleh rumus
$$f_{\mu,\sigma}(x) = \frac{1}{\sqrt {2\pi}\sigma}e^{-(x-\mu)^2/(2\sigma^2)}\ ,$$
seperti telah kita lihat dalam Bab~\ref{chp 5}.  Dalam bagian ini, kita hanya akan membahas kasus 
ketika $\mu = 0$ dan $\sigma = 1$.  Fungsi densitas normal khusus ini akan kita sebut densitas normal 
{\emx baku}, dan kita lambangkan dengan $\phi(x)$:
$$\phi(x) = \frac {1}{\sqrt{2\pi}}e^{-x^2/2}\ .$$
Grafik fungsi ini diberikan dalam Gambar~\ref{fig 9.0}.  Dapat ditunjukkan bahwa luas
di bawah setiap densitas normal sama dengan 1.
\putfig{3truein}{PSfig9-0}
{Densitas normal baku.}{fig 9.0} 
\par
Secara cukup umum, Teorema Limit Pusat menjelaskan apa yang
terjadi ketika kita menjumlahkan sejumlah besar peubah acak yang saling bebas,
yang masing-masing hanya menyumbang sedikit terhadap jumlah keseluruhan.  Dalam bagian ini kita
akan membahas penerapan teorema ini pada percobaan Bernoulli, sedangkan dalam
Bagian~\ref{sec 9.3} kita akan mempertimbangkan proses yang lebih umum.  Kita akan membahas
teorema ini untuk kasus ketika setiap peubah acak berdistribusi identik,
tetapi, dengan syarat tertentu, teorema ini tetap berlaku meskipun setiap
peubah acak memiliki distribusi yang berbeda.


\subsection*{Percobaan Bernoulli}
Perhatikan suatu proses percobaan Bernoulli dengan peluang keberhasilan~$p$ pada setiap
percobaan.  Tetapkan $X_i = 1$ atau~0 menurut apakah hasil ke-$i$ merupakan keberhasilan atau
kegagalan, dan tetapkan $S_n = X_1 + X_2 +\cdots+ X_n$.  Maka $S_n$ adalah banyaknya
keberhasilan dalam $n$ percobaan.  Kita tahu bahwa distribusi $S_n$ diberikan oleh peluang binomial
$b(n,p,j)$.  Dalam Bagian~\ref{sec 3.2}, kita telah menggambarkan
distribusi ini untuk $p = .3$ dan $p = .5$ pada berbagai nilai~$n$ (lihat Gambar~\ref{fig 3.8}).
\par
Kita melihat bahwa nilai maksimum distribusi muncul di dekat nilai harapan
$np$, sehingga grafik pakunya bergeser ke kanan ketika
$n$ bertambah.  Selain itu, nilai maksimum ini mendekati 0 ketika $n$ bertambah,
sehingga grafik paku tersebut semakin mendatar.

\subsection*{Jumlah Terstandardisasi}
Kita dapat mencegah pergeseran grafik paku ini dengan mengurangkan banyaknya
keberhasilan yang diharapkan, $np$, dari~$S_n$, sehingga diperoleh peubah acak baru $S_n -
np$.  Sekarang nilai maksimum distribusi akan selalu berada di dekat 0.
\par
Untuk mencegah grafik paku ini melebar, kita dapat menormalkan $S_n - np$ agar
memiliki varians~1 dengan membaginya dengan simpangan baku $\sqrt{npq}$ (lihat
Latihan~\ref{sec 6.2}.\ref{exer 6.2.13}~dan~Latihan~\ref{sec 6.2}.\ref{exer 6.2.17}).

\begin{definition}
\emx {Jumlah terstandardisasi}\index{jumlah terstandardisasi} dari $S_n$ diberikan oleh
\[
S_n^* = \frac {S_n - np}{\sqrt{npq}}\ .
\]
$S_n^*$ selalu memiliki nilai harapan~0 dan varians~1.
\end{definition}

Misalkan kita menggambar grafik paku dengan paku-paku yang ditempatkan pada nilai-nilai yang mungkin
dari~$S_n^*$: $x_0$,~$x_1$, \dots,~$x_n$, dengan

\begin{equation}
x_j = \frac {j - np}{\sqrt{npq}}\ .
\label{eq 9.1}
\end{equation}                                                        
Kita menetapkan tinggi paku pada $x_j$ sama dengan nilai distribusi $b(n, p, j)$.  Contoh
grafik paku terstandardisasi ini, dengan $n = 270$ dan $p = .3$, ditunjukkan dalam Gambar~\ref{fig 9.1}.
Grafik ini berbentuk lonceng yang sangat indah.  Kita ingin mencocokkan suatu densitas normal dengan
grafik paku ini.  Pilihan yang paling jelas ialah densitas normal baku karena pusatnya berada di
0, sama seperti grafik paku terstandardisasi tersebut.  Dalam gambar ini, kita telah menggambar densitas normal
baku itu.  Pembaca akan melihat bahwa sesuatu yang tidak menyenangkan telah terjadi: meskipun bentuk
kedua grafik sama, tingginya sangat berbeda.

\putfig{4truein}{PSfig9-1}
{Distribusi binomial ternormalisasi dan densitas normal baku.}{fig 9.1} 

\par
Jika kita ingin kedua grafik itu saling berimpit, salah satunya harus diubah; kita memilih mengubah
grafik paku.   Karena bentuk kedua grafik tampak cukup serupa, kita akan mencoba mengubah
grafik paku tanpa mengubah bentuknya.  Perbedaan tinggi terjadi karena jumlah
tinggi semua paku sama dengan 1, sedangkan luas di bawah densitas normal baku sama dengan 1. 
Jika kita menggambar kurva kontinu melalui puncak paku-paku tersebut dan mencari luas di bawah
kurva ini, kira-kira kita akan memperoleh jumlah tinggi paku dikalikan
dengan jarak antara dua paku berurutan, yang akan kita sebut $\epsilon$.  Karena jumlah
tinggi paku sama dengan satu, luas di bawah kurva ini kira-kira $\epsilon$. 
Jadi, untuk mengubah grafik paku agar luas di bawah kurva ini bernilai 1, kita hanya perlu
mengalikan tinggi paku dengan $1/\epsilon$.  Dari Persamaan~\ref{eq 9.1}, mudah dilihat
bahwa 
$$\epsilon = \frac {1}{\sqrt {npq}}\ .$$ 
\putfig{4truein}{PSfig9-2}
{Grafik paku yang dikoreksi bersama densitas normal baku.}{fig 9.2} 
Dalam Gambar~\ref{fig 9.2} kita menampilkan jumlah terstandardisasi $S^*_n$ untuk $n = 270$ dan $p = .3$, 
setelah tingginya dikoreksi, bersama densitas normal baku.  (Gambar ini
dihasilkan dengan program {\bf CLTBernoulliPlot}.)\index{CLTBernoulliPlot (program)}  Pembaca
akan melihat bahwa densitas normal baku sangat cocok dengan grafik paku yang tingginya telah dikoreksi. 
Sesungguhnya, salah satu versi Teorema Limit Pusat (lihat Teorema~\ref{thm 9.1.1}) menyatakan bahwa ketika $n$
bertambah, densitas normal baku akan semakin baik dalam mengaproksimasi
grafik paku dengan tinggi terkoreksi yang bersesuaian dengan proses percobaan Bernoulli dengan $n$ suku. 
\par
Tetapkan suatu nilai~$x$ pada sumbu-$x$ dan suatu bilangan bulat positif tetap $n$.  Dengan menggunakan
Persamaan~\ref{eq 9.1}, titik $x_j$ yang paling dekat dengan $x$ memiliki subskrip $j$ yang diberikan oleh
rumus
$$
j = \langle np + x \sqrt{npq} \rangle\ ,
$$
dengan $\langle a \rangle$ menyatakan bilangan bulat yang paling dekat dengan~$a$.  
Dengan demikian, tinggi
paku di atas $x_j$ adalah
$$
\sqrt{npq}\,b(n,p,j) = \sqrt{npq}\,b(n,p,\langle np + x_j \sqrt{npq}
\rangle)\ .
$$
Untuk $n$ yang besar, kita telah melihat bahwa tinggi paku sangat dekat dengan tinggi
densitas normal pada $x$.  Hal ini mengisyaratkan teorema berikut.

\begin{theorem}{\bf (Teorema Limit Pusat untuk Distribusi Binomial)}\label{thm
9.1.1}\index{Teorema Limit Pusat!untuk Distribusi Binomial}
Untuk distribusi binomial $b(n,p,j)$ kita memiliki
$$
\lim_{n \to \infty} \sqrt{npq}\,b(n,p,\langle np + x\sqrt{npq} \rangle) = \phi(x)\ ,
$$
dengan $\phi(x)$ adalah densitas normal baku.
\par
Bukti teorema ini dapat diperoleh dengan menggunakan aproksimasi Stirling dari
Bagian~\ref{sec 3.1}.  Kita mengilustrasikan metode pembuktian ini dengan mempertimbangkan
kasus $x = 0$.  Dalam kasus ini, teorema tersebut menyatakan bahwa
$$
\lim_{n \to \infty} \sqrt{npq}\,b(n,p,\langle np \rangle) = \frac 1{\sqrt{2\pi}}
= .3989\ldots\ .
$$
Untuk menyederhanakan perhitungan, kita asumsikan bahwa $np$ merupakan bilangan bulat, sehingga
$\langle np \rangle = np$.  Maka
$$
\sqrt{npq}\,b(n,p,np) = \sqrt{npq}\,p^{np}q^{nq} \frac {n!}{(np)!\,(nq)!}\ .
$$
Ingat bahwa rumus Stirling (lihat Teorema~\ref{thm 3.3}) menyatakan bahwa
$$
n! \sim \sqrt{2\pi n}\,n^n e^{-n} \qquad \mbox {ketika \,\,\,} n \to \infty\ .
$$
Dengan menggunakan hasil ini, kita memperoleh
$$
\sqrt{npq}\,b(n,p,np) \sim \frac {\sqrt{npq}\,p^{np}q^{nq} \sqrt{2\pi n}\,n^n
e^{-n}}{\sqrt{2\pi np} \sqrt{2\pi nq}\,(np)^{np} (nq)^{nq} e^{-np} e^{-nq}}\ ,
$$
yang menyederhana menjadi $1/\sqrt{2\pi}$.
\end{theorem}
\subsection*{Mengaproksimasi Distribusi Binomial}\index{distribusi binomial!aproksimasi}
Kita dapat menggunakan Teorema~\ref{thm 9.1.1} untuk mencari aproksimasi nilai fungsi
distribusi binomial.  Jika kita ingin mencari aproksimasi untuk $b(n, p, j)$, kita tetapkan
$$j = np + x\sqrt{npq}$$
dan menyelesaikannya untuk $x$, sehingga diperoleh
$$x = {{j-np}\over{\sqrt{npq}}}\ .$$
Teorema~\ref{thm 9.1.1} kemudian menyatakan bahwa 
$$\sqrt{npq}\,b(n,p,j)$$
kira-kira sama dengan $\phi(x)$, sehingga
\begin{eqnarray*}
b(n,p,j) &\approx& {{\phi(x)}\over{\sqrt{npq}}}\\
&=& {1\over{\sqrt{npq}}} \phi\biggl({{j-np}\over{\sqrt{npq}}}\biggr)\ .\\
\end{eqnarray*}


\begin{example}\label{exam 9.1}
Mari kita perkirakan peluang muncul tepat 55 sisi kepala dalam 100 lemparan koin. 
Untuk kasus ini $np = 100 \cdot 1/2 = 50$ dan $\sqrt{npq} = \sqrt{100 \cdot 1/2
\cdot 1/2} = 5$.  Jadi, $x_{55} = (55 - 50)/5 = 1$ dan

\begin{eqnarray*}
P(S_{100} = 55) \sim \frac {\phi(1)}5 &=& \frac 15 \left( \frac 1{\sqrt{2\pi}}e^{-1/2} \right) \\
                                   &=& .0484\ . \\
\end{eqnarray*}

Hingga empat tempat desimal, nilai sebenarnya adalah .0485, sehingga
aproksimasi tersebut sangat baik.
\end{example}

Program {\bf CLTBernoulliLocal}\index{CLTBernoulliLocal (program)} mengilustrasikan
aproksimasi ini untuk setiap pilihan $n$,~$p$, dan~$j$.  Kita telah menjalankan program ini untuk dua
contoh.  Yang pertama adalah peluang muncul tepat 50~sisi kepala dalam 100~lemparan koin;
perkiraannya .0798, sedangkan nilai sebenarnya hingga empat tempat desimal ialah .0796.  Contoh
kedua adalah peluang muncul mata dadu enam tepat delapan kali dalam 36~lemparan dadu; di sini perkiraannya
.1093, sedangkan nilai sebenarnya hingga empat tempat desimal ialah .1196.
\par
Setiap peluang binomial menuju~0 ketika $n$ menuju tak hingga.  Dalam
kebanyakan penerapan, kita tidak tertarik pada peluang terjadinya suatu hasil
tertentu, melainkan pada peluang bahwa hasil tersebut berada dalam suatu
interval, misalnya interval $[a, b]$.  Untuk mencari peluang ini, kita menjumlahkan
tinggi grafik paku untuk nilai~$j$ antara $a$~dan~$b$.  Ini sama dengan
menanyakan peluang bahwa jumlah terstandardisasi $S_n^*$ berada antara $a^*$ dan 
$b^*$, dengan $a^*$ dan $b^*$ adalah nilai terstandardisasi dari $a$ dan $b$.  Namun, ketika $n$
menuju tak hingga, jumlah luas ini dapat diharapkan mendekati luas
di bawah densitas normal baku antara $a^*$~dan~$b^*$.  \emx {Teorema Limit Pusat}
menyatakan bahwa hal ini memang terjadi.

\begin{theorem}{\bf (Teorema Limit Pusat untuk Percobaan Bernoulli)}\index{Teorema Limit
Pusat!untuk Percobaan Bernoulli} 
Misalkan $S_n$ banyaknya keberhasilan dalam $n$ percobaan Bernoulli dengan peluang keberhasilan $p$, dan misalkan
$a$ dan $b$ dua bilangan real tetap.  Maka
$$\lim_{n \rightarrow \infty} P\biggl(a \le \frac{S_n - np}{\sqrt{npq}} \le b\biggr) = \int_a^b \phi(x)\,dx\ .$$
\end{theorem}

Teorema ini dapat dibuktikan dengan menjumlahkan aproksimasi terhadap $b(n,p,k)$ yang diberikan dalam 
Teorema~\ref{thm 9.1.1}.\choice{\footnote{Teorema ini juga merupakan kasus khusus dari Teorema Limit Pusat yang lebih
umum. Lihat Bagian 10.3 dalam buku Grinstead--Snell lengkap.}}{Teorema ini juga merupakan kasus khusus dari Teorema Limit Pusat yang lebih
umum (lihat Bagian~\ref{sec 10.3}).}
\par
Kita tahu dari kalkulus bahwa integral di ruas kanan persamaan ini
sama dengan luas di bawah grafik densitas normal baku $\phi(x)$ antara
$a$~dan~$b$.  Luas ini kita lambangkan dengan $\NA(a^*, b^*)$.  Sayangnya, tidak ada cara sederhana untuk
mengintegralkan fungsi
$e^{-x^2/2}$, sehingga kita harus menggunakan tabel nilai atau program
integrasi numerik.  (Lihat Gambar~\ref{tabl 9.1} untuk nilai $\NA(0, z)$.  Tabel yang lebih lengkap 
diberikan dalam Lampiran~A.)
\putfig{4.5truein}{PSfig9-2-5}
{Tabel nilai $\NA(0,z)$, yaitu luas normal dari 0~hingga~$z$.}{tabl 9.1}
\par
Dari simetri densitas normal baku, jelas bahwa luas seperti luas
antara $-2$~dan~3 dapat dicari dari tabel ini dengan menjumlahkan luas dari 0~hingga~2
(sama dengan luas dari $-2$~hingga~0) dan luas dari 0~hingga~3.
\par

\subsection*{Aproksimasi Peluang Binomial}

Misalkan $S_n$ berdistribusi binomial dengan parameter $n$ dan $p$.  Kita telah melihat bahwa
teorema di atas menunjukkan cara memperkirakan peluang berbentuk 
\begin{equation}
P(i \le S_n \le j)\ ,
\label{eq 9.2}
\end{equation}
dengan $i$ dan $j$ bilangan bulat antara 0 dan $n$.  Seperti telah kita lihat, distribusi binomial dapat
direpresentasikan sebagai grafik paku, dengan paku pada bilangan bulat antara 0 dan $n$, serta
tinggi paku ke-$k$ diberikan oleh $b(n, p, k)$.  Untuk nilai $n$ berukuran sedang, jika kita
menstandardisasi grafik paku ini dan mengubah tinggi paku-pakunya dengan cara yang
dijelaskan di atas, jumlah tinggi paku diaproksimasi oleh luas di bawah
densitas normal baku antara $i^*$ dan $j^*$.  Ternyata, aproksimasi yang sedikit lebih akurat
diberikan oleh luas di bawah densitas normal baku antara nilai-nilai terstandardisasi
yang bersesuaian dengan $(i - 1/2)$ dan $(j + 1/2)$; nilai-nilai ini adalah
$$i^* = \frac{i - 1/2 - np}{\sqrt {npq}}$$
dan
$$j^* = \frac{j + 1/2 - np}{\sqrt {npq}}\ .$$
Dengan demikian,
$$P(i \le S_n \le j) \approx \NA\Biggl({{i - {1\over 2} - np}\over{\sqrt {npq}}} ,
{{j + {1\over 2} - np}\over{\sqrt {npq}}}\Biggr)\ .$$
Perlu ditekankan bahwa hasil yang diperoleh dengan menggunakan Teorema Limit Pusat hanyalah
aproksimasi dan terkadang tidak terlalu dekat dengan nilai sebenarnya (lihat Latihan~\ref{exer 9.2.111}). 
\par
Sekarang kita ilustrasikan gagasan ini dengan beberapa contoh.

\begin{example}\label{exam 9.2}
Sebuah koin dilempar 100 kali.  Perkirakan peluang bahwa banyaknya kemunculan sisi kepala
berada antara 40~dan~60 (dalam matematika, kata ``antara" mencakup kedua titik ujung).  Banyaknya
kemunculan sisi kepala yang diharapkan adalah
$100
\cdot 1/2 = 50$, dan simpangan baku banyaknya kemunculan sisi kepala adalah $\sqrt{100 \cdot 1/2
\cdot 1/2} = 5$.  Jadi, karena $n = 100$ cukup besar, kita memperoleh
\begin{eqnarray*}
P(40 \le S_n \le 60) &\approx& 
P\left( \frac {39.5 - 50}5 \le S_n^* \le \frac {60.5 - 50}5 \right) \\
                     &=& P(-2.1 \le S_n^* \le 2.1) \\
                     &\approx& \NA(-2.1,2.1) \\ 
                     &=& 2\NA(0,2.1) \\
                     &\approx& .9642\ .  
\end{eqnarray*}
Nilai sebenarnya hingga lima tempat desimal adalah .96480.
\par
Perhatikan bahwa dalam kasus ini kita menanyakan peluang bahwa hasilnya
tidak menyimpang lebih dari dua simpangan baku dari nilai harapan.  Seandainya
kita menanyakan peluang bahwa banyaknya keberhasilan berada antara 35~dan~65,
rentang ini akan mewakili tiga simpangan baku dari rata-rata dan, dengan menggunakan koreksi 1/2,
perkiraan kita adalah luas di bawah kurva normal baku antara $-3.1$~dan~3.1, 
atau $2\NA(0,3.1) = .9980$.  Jawaban sebenarnya dalam kasus ini, hingga lima tempat, adalah .99821.  
\end{example}


Penting untuk mengerjakan beberapa soal secara manual agar memahami pengubahan
dari suatu pertidaksamaan ke pertidaksamaan yang berkaitan dengan peubah
terstandardisasi.  Setelah itu, kita dapat menggunakan program komputer yang melakukan
pengubahan ini, termasuk koreksi 1/2.  Program {\bf
CLTBernoulliGlobal}\index{CLTBernoulliGlobal} merupakan program semacam itu untuk memperkirakan peluang
berbentuk
$P(a \leq S_n \leq b)$.

\begin{example}\label{exam 9.3}
Dartmouth College ingin memiliki 1050 mahasiswa tahun pertama.  Perguruan tinggi ini tidak dapat
menampung lebih dari 1060.  Asumsikan bahwa setiap pelamar menerima tawaran dengan
peluang~.6 dan penerimaan tawaran dapat dimodelkan dengan percobaan Bernoulli. 
Jika perguruan tinggi tersebut menawarkan tempat kepada 1700 pelamar, berapa peluang bahwa terlalu
banyak dari mereka akan menerima tawaran tersebut?
\par
Jika menawarkan tempat kepada 1700 pelamar, banyaknya mahasiswa yang diharapkan mendaftar ulang adalah
$.6 \cdot 1700 = 1020$.  Simpangan baku banyaknya pelamar yang menerima tawaran adalah
$\sqrt{1700 \cdot .6 \cdot .4} \approx 20$.  Jadi, kita ingin memperkirakan
peluang
\begin{eqnarray*}
P(S_{1700} > 1060) &=& P(S_{1700} \ge 1061) \\
&=& P\left( S_{1700}^* \ge \frac {1060.5 - 1020}{20} \right) \\
                   &=& P(S_{1700}^* \ge 2.025)\ .
\end{eqnarray*}

Dari Tabel~\ref{tabl 9.1}, jika kita melakukan interpolasi, peluang ini diperkirakan
sebesar $.5 - .4784 = .0216$.  Jadi, perguruan tinggi tersebut cukup aman menggunakan
kebijakan penerimaan ini.
\end{example}


\subsection*{Penerapan pada Statistika}\index{statistika!penerapan Teorema Limit
Pusat}
Ada banyak pertanyaan penting dalam bidang statistika yang dapat dijawab dengan menggunakan
Teorema Limit Pusat untuk proses percobaan yang saling bebas.  Contoh berikut cukup
sering dijumpai dalam berita.  Contoh lain penerapan Teorema
Limit Pusat pada statistika diberikan dalam Bagian~\ref{sec 9.3}.

\begin{example}\label{exam 9.4.1}
Kita sering membaca bahwa suatu jajak pendapat\index{jajak pendapat} dilakukan untuk memperkirakan proporsi orang
dalam populasi tertentu yang lebih menyukai satu calon daripada calon lain dalam pemilihan dengan dua calon. 
(Model ini juga berlaku pada pemilihan dengan lebih dari dua calon ketika kita membandingkan calon $A$ dan $B$, serta pada dua pilihan
dalam suatu pemungutan suara.)  Tentu saja, penyelenggara jajak pendapat tidak mungkin menanyakan pilihan setiap orang. 
Sebagai gantinya, dipilih suatu subhimpunan populasi yang disebut sampel\index{sampel}, lalu
setiap orang dalam sampel ditanya tentang pilihannya.  Misalkan $p$ proporsi sebenarnya dari orang dalam
populasi yang mendukung calon $A$, dan misalkan $q = 1-p$.  Jika kita memilih sampel berukuran $n$
dari populasi, pilihan orang-orang dalam sampel dapat direpresentasikan oleh peubah acak
$X_1,\ X_2,\ \ldots,\ X_n$, dengan $X_i = 1$ jika orang ke-$i$ mendukung calon $A$,
dan $X_i = 0$ jika orang ke-$i$ mendukung calon $B$.  Misalkan $S_n = X_1 + X_2 + \cdots + X_n$. 
Jika setiap subhimpunan berukuran $n$ dipilih dengan peluang yang sama, maka $S_n$ berdistribusi
hipergeometrik.  Jika $n$ kecil dibandingkan ukuran populasi (seperti yang biasanya terjadi dalam
praktik), maka $S_n$ kira-kira berdistribusi binomial dengan parameter $n$ dan $p$.
\par
Penyelenggara jajak pendapat ingin memperkirakan nilai $p$.  Perkiraan untuk $p$ diberikan oleh nilai
$\bar p = S_n/n$, yaitu proporsi orang dalam sampel yang mendukung calon $B$.
Teorema Limit Pusat menyatakan bahwa peubah acak $\bar p$ kira-kira berdistribusi
normal.  (Sesungguhnya, versi Teorema Limit Pusat kita menyatakan bahwa fungsi
distribusi peubah acak
$$S_n^* = \frac{S_n - np}{\sqrt{npq}}$$
diaproksimasi oleh densitas normal baku.)  Namun, kita memiliki
$$\bar p = \frac{S_n - np}{\sqrt {npq}}\sqrt{\frac{pq}{n}}+p\ ,$$
yaitu, $\bar p$ hanyalah fungsi linear dari $S_n^*$.  Karena distribusi $S_n^*$
diaproksimasi oleh densitas normal baku, distribusi peubah acak $\bar p$
juga harus berbentuk lonceng.  Kita juga tahu cara menyatakan rata-rata dan simpangan baku $\bar p$
dalam
$p$ dan
$n$.  Rata-rata $\bar p$ adalah $p$, dan simpangan bakunya ialah
$$\sqrt{\frac{pq}{n}}\ .$$
Jadi, mudah untuk menuliskan versi terstandardisasi dari $\bar p$, yaitu
$$\bar p^* = \frac{\bar p - p}{\sqrt{pq/n}}\ .$$
\par
Karena distribusi versi terstandardisasi dari $\bar p$ diaproksimasi oleh densitas
normal baku, kita tahu, misalnya, bahwa 95\% nilainya akan berada dalam jarak dua simpangan baku
dari rata-ratanya, dan hal yang sama berlaku untuk $\bar p$.  Jadi, kita memiliki
$$P\left(p - 2\sqrt{\frac{pq}{n}} < \bar p < p + 2\sqrt{\frac{pq}{n}}\right) \approx .954\
.$$
Penyelenggara jajak pendapat tidak mengetahui $p$ atau $q$, tetapi dapat menggantikannya dengan $\bar p$ dan $\bar q = 1 -
\bar p$ tanpa risiko yang berarti.  Dengan gagasan ini, pernyataan di atas
setara dengan pernyataan
$$P\left(\bar p - 2\sqrt{\frac{\bar p \bar q}{n}} < p <
\bar p + 2\sqrt{\frac{\bar p \bar q}{n}}\right) \approx .954\ .$$
Interval yang dihasilkan,
$$
\left( \bar p - \frac {2\sqrt{\bar p \bar q}}{\sqrt n},\ 
\bar p + \frac {2\sqrt{\bar p \bar q}}{\sqrt n} \right)
$$
disebut \emx {interval kepercayaan 95 persen}\index{interval kepercayaan} untuk nilai
$p$ yang tidak diketahui.  Nama ini berasal dari kenyataan bahwa jika kita menggunakan metode ini untuk
memperkirakan $p$ pada sejumlah besar sampel, kita dapat mengharapkan bahwa pada sekitar
95~persen sampel, nilai sebenarnya dari~$p$ termuat dalam interval kepercayaan
yang diperoleh dari sampel.  Dalam Latihan~\ref{exer 9.1.11}, Anda diminta
menulis program untuk mengilustrasikan bahwa hal ini memang terjadi.
\par
Penyelenggara jajak pendapat dapat menentukan nilai $n$.  Jadi, jika ingin membuat interval
kepercayaan 95\% dengan panjang 6\%, nilai $n$ harus dipilih sedemikian sehingga
$$\frac {2\sqrt{\bar p \bar q}}{\sqrt n} \le .03\ .$$
Dengan menggunakan fakta bahwa $\bar p \bar q \le 1/4$ berapa pun nilai $\bar p$,
mudah ditunjukkan bahwa jika nilai $n$ dipilih sedemikian sehingga
$$\frac{1}{\sqrt n} \le .03\ ,$$
batas tersebut akan terpenuhi.  Ini setara dengan memilih
$$n \ge 1111\ .$$
Jadi, jika penyelenggara jajak pendapat memilih, misalnya, $n$ sebesar 1200 dan menghitung $\bar p$ menggunakan
sampel berukuran 1200, maka 19 dari 20 kali (yaitu 95\% dari seluruh percobaan), interval kepercayaannya,
yang panjangnya 6\%, akan memuat nilai sebenarnya dari $p$.  Jenis interval kepercayaan ini
biasanya dilaporkan dalam berita sebagai berikut: survei ini memiliki margin galat 3\%.\index{margin
galat}  Sesungguhnya, sebagian besar survei yang diberitakan di surat kabar memiliki ukuran sampel
sekitar 1000.  Fakta yang agak mengejutkan ialah bahwa ukuran populasi tampaknya
tidak memengaruhi ukuran sampel yang diperlukan untuk memperoleh interval kepercayaan 95\% bagi $p$ dengan
margin galat tertentu.  Untuk melihatnya, perhatikan bahwa nilai $n$ yang diperlukan hanya bergantung pada
bilangan .03, yaitu margin galat.  Dengan kata lain, baik ukuran populasi
100{,}000 maupun 100{,}000{,}000, penyelenggara jajak pendapat hanya perlu memilih sampel berukuran sekitar 1200 untuk
memperoleh ketelitian perkiraan $p$ yang sama.  (Kita memang menggunakan fakta bahwa ukuran sampel kecil
dibandingkan ukuran populasi ketika menyatakan bahwa $S_n$ kira-kira berdistribusi
binomial.)
\par
Dalam Gambar~\ref{fig 9.2.1}, kita menampilkan hasil simulasi proses jajak pendapat.  Ukuran populasi
adalah 100{,}000, dan untuk populasi tersebut, $p = .54$.  Ukuran sampel dipilih sebesar
1200.  Grafik paku menunjukkan distribusi $\bar p$ untuk 10{,}000 sampel yang dipilih secara
acak.   Dalam simulasi ini, program mencatat banyaknya sampel dengan
$\bar p$ berada dalam jarak 3\% dari .54.  Jumlahnya 9648, yang mendekati
95\% dari banyaknya sampel yang digunakan.
\putfig{4truein}{PSfig9-2-1}{Simulasi jajak pendapat.}{fig 9.2.1} 
\par                        
Cara lain untuk memahami makna interval kepercayaan ditunjukkan dalam Gambar~\ref{fig
9.2.2}.  Dalam gambar ini, kita menampilkan 100 interval kepercayaan yang diperoleh dengan menghitung $\bar p$
untuk 100 sampel berbeda berukuran 1200 dari populasi yang sama seperti sebelumnya.  Pembaca dapat melihat
bahwa sebagian besar interval kepercayaan ini (tepatnya 96) memuat nilai sebenarnya dari $p$.
\putfig{4.5truein}{PSfig9-2-2}{Simulasi interval kepercayaan.}{fig 9.2.2} 
\par
Jajak Pendapat Gallup\index{Jajak Pendapat Gallup} telah menggunakan teknik jajak pendapat ini dalam setiap
pemilihan Presiden\index{pemilihan Presiden} sejak 1936 (serta dalam tak terhitung banyaknya
pemilihan lain).  Tabel~\ref{table 9.1}\footnote{The Gallup Poll Monthly, November 1992, 
No.\ 326, p.\ 33.  Dilengkapi dengan bantuan Lydia K. Saab, The Gallup Organization.} menunjukkan
hasil upaya tersebut.  Pembaca akan melihat bahwa sebagian besar aproksimasi terhadap $p$ berada
dalam jarak 3\% dari nilai sebenarnya dari $p$.  Ukuran sampel untuk jajak pendapat ini biasanya sekitar
1500.  (Dalam tabel, persentase prediksi maupun persentase sebenarnya untuk calon pemenang mengacu
pada persentase suara di antara partai politik ``utama". Dalam kebanyakan pemilihan terdapat
dua partai utama, tetapi dalam beberapa pemilihan terdapat tiga.) 
\begin{table}
\centering
\begin{tabular}{clccc} Tahun &$\,$ Calon &Survei Akhir &Hasil 
&Deviasi \\
&Pemenang&Gallup&Pemilu&\\
\hline
1936 & Roosevelt  & 55.7\% & 62.5\% & 6.8\%\\
1940 & Roosevelt  & 52.0\% & 55.0\% & 3.0\%\\
1944 & Roosevelt  & 51.5\% & 53.3\% & 1.8\%\\
1948 & Truman     & 44.5\% & 49.9\% & 5.4\%\\
1952 & Eisenhower & 51.0\% & 55.4\% & 4.4\%\\
1956 & Eisenhower & 59.5\% & 57.8\% & 1.7\%\\
1960 & Kennedy    & 51.0\% & 50.1\% & 0.9\%\\
1964 & Johnson    & 64.0\% & 61.3\% & 2.7\%\\
1968 & Nixon      & 43.0\% & 43.5\% & 0.5\%\\
1972 & Nixon      & 62.0\% & 61.8\% & 0.2\%\\
1976 & Carter     & 48.0\% & 50.0\% & 2.0\%\\
1980 & Reagan     & 47.0\% & 50.8\% & 3.8\%\\
1984 & Reagan     & 59.0\% & 59.1\% & 0.1\%\\
1988 & Bush       & 56.0\% & 53.9\% & 2.1\%\\
1992 & Clinton    & 49.0\% & 43.2\% & 5.8\%\\
1996 & Clinton    & 52.0\% & 50.1\% & 1.9\%\\
\end{tabular}
\caption{Catatan akurasi Jajak Pendapat Gallup.}
\label{table 9.1}
\end{table}
\par
Teknik ini juga berperan penting dalam mengevaluasi efektivitas obat dalam
bidang kedokteran.  Sebagai contoh, terkadang kita ingin mengetahui berapa proporsi pasien
yang akan terbantu oleh obat baru.  Proporsi ini dapat diperkirakan dengan memberikan obat kepada suatu subhimpunan
pasien, lalu menentukan proporsi anggota sampel yang terbantu oleh obat tersebut.
\end{example}



\subsection*{Catatan Sejarah}
Teorema Limit Pusat untuk percobaan Bernoulli pertama kali dibuktikan oleh Abraham\linebreak[4]
de~Moivre\index{de MOIVRE, A.} dan dimuat dalam bukunya, \emx {The Doctrine of Chances,} yang pertama kali
diterbitkan pada 1718.\footnote{A. de Moivre, \emx {The Doctrine of Chances,} edisi~ke-3.\ (London: Millar,
1756).}

De Moivre menjalani usia 18~hingga~21 tahun di penjara di Prancis karena latar
belakang Protestannya.  Setelah dibebaskan, ia meninggalkan Prancis menuju Inggris, tempat
ia bekerja sebagai guru privat bagi putra-putra bangsawan.  Newton pernah memberikan satu eksemplar
\emx {Principia Mathematica} karyanya kepada Earl of Devonshire.  Menurut kisah yang beredar,
ketika de~Moivre mengajar di kediaman sang Earl, ia menemukan karya Newton
dan merasa bahwa karya itu melampaui pemahamannya.  Konon, ia kemudian membeli eksemplarnya
sendiri dan merobeknya menjadi halaman-halaman terpisah, mempelajarinya halaman demi halaman sambil berjalan
mengelilingi London menuju tempat ia mengajar.  De~Moivre sering mengunjungi kedai-kedai kopi di
London, tempat ia memulai karyanya dalam probabilitas dengan menghitung peluang bagi para
penjudi.  Ia juga bertemu Newton di salah satu kedai kopi tersebut dan keduanya menjadi
sahabat dekat.  De~Moivre mendedikasikan bukunya kepada Newton.

\emx {The Doctrine of Chances} menyajikan teknik untuk menyelesaikan beragam
masalah perjudian.  Di tengah pembahasan masalah-masalah perjudian ini,
de~Moivre dengan cukup rendah hati memperkenalkan buktinya atas Teorema Limit Pusat,
dengan menulis
\begin{quote}
Suatu Metode untuk mengaproksimasi Jumlah Suku-suku Binomial $(a + b)^n$
yang dikembangkan menjadi Deret, yang darinya diturunkan beberapa Kaidah praktis untuk
memperkirakan Derajat Keyakinan yang patut diberikan kepada
Percobaan.\footnote{ibid., p.~243.}
\end{quote}
Bukti de Moivre menggunakan aproksimasi faktorial yang sekarang kita sebut
rumus Stirling.  De~Moivre menyatakan bahwa ia telah memperoleh rumus ini sebelum
Stirling, tetapi tanpa menentukan nilai tepat konstanta
$\sqrt{2\pi}$.  Meskipun mengatakan bahwa nilai tepat ini sebenarnya tidak perlu
diketahui, ia mengakui bahwa pengetahuan tentangnya ``memberikan Keanggunan yang istimewa pada
Penyelesaian."
\par
Bukti lengkap dan pembahasan menarik tentang kehidupan 
de~Moivre dapat ditemukan dalam buku \emx {Games, Gods and Gambling} karya F.~N.
David.\index{DAVID, F. N.}\footnote{F. N. David, \emx {Games, Gods and Gambling} (London:
Griffin, 1962).}

\pagebreak[4]

\exercises
\begin{LJSItem}

\i\label{exer 9.1.1}  Misalkan $S_{100}$ menyatakan banyaknya kemunculan sisi kepala dalam 100 lemparan
koin seimbang.  Gunakan Teorema Limit Pusat untuk memperkirakan
\begin{enumerate}

\item  $P(S_{100} \leq 45)$.

\item  $P(45 < S_{100} < 55)$.

\item  $P(S_{100} > 63)$.

\item $P(S_{100} < 57)$.
\end{enumerate}

\i\label{exer 9.1.2}  Misalkan $S_{200}$ menyatakan banyaknya kemunculan sisi kepala dalam 200 lemparan
koin seimbang.  Perkirakan
\begin{enumerate}

\item  $P(S_{200} = 100)$. 

\item  $P(S_{200} = 90)$. 

\item  $P(S_{200} = 80)$. 
\end{enumerate}

\i\label{exer 9.1.3}  Suatu ujian benar-salah memiliki 48 soal.  June memiliki peluang~3/4 untuk
menjawab suatu soal dengan benar.  April hanya menebak pada setiap soal.  Nilai
lulus diperoleh dengan 30~jawaban benar atau lebih.  Bandingkan peluang June
lulus ujian dengan peluang April lulus.

\i\label{exer 9.1.4}  Misalkan $S$ menyatakan banyaknya kemunculan sisi kepala dalam 1{,}000{,}000 lemparan
koin seimbang.  Gunakan (a)~ketaksamaan Chebyshev dan (b)~Teorema Limit Pusat untuk
memperkirakan peluang bahwa $S$ berada antara 499{,}500 dan 500{,}500.  Gunakan
dua metode yang sama untuk memperkirakan peluang bahwa $S$ berada antara
499{,}000 dan 501{,}000, serta peluang bahwa $S$ berada antara 498{,}500
dan 501{,}500.

\i\label{exer 9.1.5}  Seorang pemain pemula direkrut oleh klub bisbol dengan asumsi bahwa ia akan
memiliki rata-rata pukulan .300.  (Rata-rata pukulan adalah rasio banyaknya
pukulan berhasil terhadap banyaknya kesempatan memukul.)  Pada tahun pertama, ia mendapat 300
kesempatan memukul dan rata-rata pukulannya .267.  Asumsikan bahwa setiap kesempatan memukul dapat dianggap sebagai
percobaan Bernoulli dengan peluang keberhasilan .3.  Apakah rata-rata serendah itu dapat
dianggap sekadar nasib buruk, atau haruskah ia dikembalikan ke liga minor? 
Berikan komentar tentang asumsi percobaan Bernoulli dalam situasi ini.

\i\label{exer 9.1.6}  Pada suatu masa, dua kereta api bersaing untuk mengangkut
1000 penumpang yang berangkat dari Chicago pada jam yang sama menuju
Los Angeles.  Asumsikan bahwa setiap penumpang memiliki peluang yang sama untuk memilih salah satu dari kedua
kereta.  Berapa banyak kursi yang harus dimiliki sebuah kereta agar peluang tersedianya kursi
bagi setiap penumpang setidaknya .99?

\i\label{exer 9.1.7}  Asumsikan bahwa, seperti dalam Contoh~\ref{exam 9.3}, Dartmouth menawarkan tempat kepada 1750
calon mahasiswa.  Berapa peluang bahwa terlalu banyak di antara mereka menerima tawaran tersebut?

\i\label{exer 9.1.8}  Sebuah klub hanya menyajikan makan malam kepada anggotanya.  Mereka duduk di meja berkapasitas
12 orang.  Dari pengamatan jangka panjang, manajer mendapati bahwa selama 95~persen
waktu terdapat antara enam dan sembilan meja penuh anggota, sedangkan pada waktu selebihnya
banyaknya meja penuh berpeluang sama untuk berada di atas atau di bawah rentang ini. 
Asumsikan bahwa setiap anggota memutuskan untuk datang dengan peluang tertentu~$p$, dan bahwa
keputusan mereka saling bebas.  Berapa banyak anggota klub tersebut?  Berapakah $p$?

\i\label{exer 9.1.9}  Misalkan $S_n$ banyaknya keberhasilan dalam $n$ percobaan Bernoulli dengan
peluang keberhasilan~.8 pada setiap percobaan.  Misalkan $A_n = S_n/n$ rata-rata keberhasilan. 
Untuk setiap kasus, berikan nilai limit dan alasan bagi jawaban Anda.
\begin{enumerate}
\item  $\lim_{n \to \infty} P(A_n = .8)$.

\item  $\lim_{n \to \infty} P(.7n < S_n < .9n)$.

\item  $\lim_{n \to \infty} P(S_n < .8n + .8\sqrt n)$.

\item  $\lim_{n \to \infty} P(.79 < A_n < .81)$.
\end{enumerate}

\i\label{exer 9.1.10}  Tentukan peluang bahwa di antara 10{,}000 digit acak, digit~3
muncul tidak lebih dari 931 kali.

\i\label{exer 9.1.11} Tulis program komputer untuk menyimulasikan 10{,}000
percobaan Bernoulli dengan peluang keberhasilan~.3 pada setiap percobaan.  Minta
program menghitung interval kepercayaan 95~persen untuk peluang
keberhasilan berdasarkan proporsi keberhasilan.  Ulangi percobaan 100 kali
dan amati berapa kali nilai sebenarnya,~.3, termuat dalam batas-batas
kepercayaan.

\i\label{exer 9.1.12}  Sebuah koin seimbang dilempar 400 kali.  Tentukan bilangan~$x$ sedemikian sehingga
peluang bahwa banyaknya kemunculan sisi kepala berada antara $200 - x$ dan $200 + x$
kira-kira .80.

\i\label{exer 9.1.13}  Mesin pembuat mi di pabrik spageti milik Spumoni menghasilkan sekitar 5~persen
mi cacat bahkan ketika disetel dengan benar.  Mi tersebut kemudian dikemas dalam
peti yang masing-masing berisi 1900 mi.  Sebuah peti diperiksa dan ternyata berisi
115 mi cacat.  Berapa perkiraan peluang menemukan setidaknya
sebanyak ini mi cacat jika mesin disetel dengan benar?

\i\label{exer 9.1.14}  Sebuah restoran melayani 400 pelanggan per hari.  Rata-rata, 20~persen
pelanggan memesan pai apel.
\begin{enumerate}
\item  Berikan rentang (yang disebut interval kepercayaan 95~persen) untuk banyaknya
potong pai apel yang dipesan pada suatu hari sehingga Anda dapat yakin 95~persen
bahwa jumlah sebenarnya akan berada dalam rentang tersebut.

\item  Berapa rata-rata banyaknya pelanggan yang harus dilayani restoran agar
setidaknya yakin 95~persen bahwa persentase pelanggan yang memesan pai pada hari itu
berada dalam rentang 19~hingga~21 persen?
\end{enumerate}

\i\label{exer 9.1.15}  Ingat bahwa jika $X$ peubah acak, \emx {fungsi distribusi
kumulatif} dari~$X$ adalah fungsi $F(x)$ yang didefinisikan oleh
$$
F(x) = P(X \leq x)\ .
$$
\begin{enumerate}
\item  Misalkan $S_n$ banyaknya keberhasilan dalam $n$ percobaan Bernoulli dengan
peluang keberhasilan~$p$.  Tulis program untuk menggambar distribusi kumulatif
untuk~$S_n$.

\item  Ubah program Anda dalam (a) untuk menggambar
distribusi kumulatif $F_n^*(x)$ dari peubah acak terstandardisasi
$$
S_n^* = \frac {S_n - np}{\sqrt{npq}}\ .
$$

\item  Definisikan \emx {distribusi normal} $N(x)$ sebagai luas di bawah
kurva normal hingga nilai~$x$.  Ubah program Anda dalam
(b) agar turut menggambar distribusi normal, lalu bandingkan
dengan distribusi kumulatif~$S_n^*$.  Lakukan ini untuk $n = 10, 50$, dan $100$.
\end{enumerate}

\i\label{exer 9.1.16}  Dalam Contoh~\ref{exam 3.12}, kita ingin menguji
hipotesis bahwa suatu bentuk aspirin baru efektif dalam 80~persen kasus,
bukan 60~persen seperti yang dilaporkan untuk aspirin standar.  Aspirin
baru diberikan kepada $n$ orang.  Jika efektif dalam $m$ kasus atau lebih,
kita menerima klaim bahwa obat baru tersebut efektif dalam 80~persen kasus;
jika tidak, kita menolak klaim itu.  Dengan menggunakan Teorema Limit Pusat, tunjukkan bahwa Anda dapat
memilih banyaknya percobaan~$n$ dan nilai kritis~$m$ sedemikian sehingga
peluang kita menolak hipotesis ketika hipotesis itu benar kurang dari .01 dan
peluang kita menerimanya ketika hipotesis itu salah juga kurang dari .01.  Tentukan
nilai terkecil~$n$ yang mencukupi untuk ini.

\i\label{exer 9.1.17}  Dalam suatu jajak pendapat, diasumsikan bahwa proporsi yang tidak diketahui, $p$,
dari masyarakat mendukung usulan undang-undang baru dan proporsi $1-p$ menentangnya.
Sampel beranggotakan $n$ orang diambil untuk memperoleh pendapat mereka.  Proporsi ${\bar
p}$ yang mendukung dalam sampel digunakan sebagai perkiraan $p$.  Dengan menggunakan Teorema Limit
Pusat, tentukan ukuran sampel yang menjamin bahwa perkiraan tersebut, dengan
peluang .95, menyimpang tidak lebih dari .01 dari nilai sebenarnya.

\i\label{exer 9.1.18} 
Sebuah uraian tentang jajak pendapat dalam suatu surat kabar menyatakan bahwa seseorang dapat 
yakin 95\% bahwa galat akibat pengambilan sampel tidak akan lebih besar 
dari plus atau minus 3 poin persentase.  Uraian sebuah jajak pendapat
\emx{New York Times}\index{New York Times} di Iowa dapat diringkas sebagai berikut:
dalam 19 dari 20 sampel semacam itu, hasil sampel diperkirakan berada dalam jarak
3 poin persentase dari hasil sensus seluruh penduduk dewasa Iowa. Ringkasan ini tidak
menerjemahkan atau mereproduksi ungkapan surat kabar. Keduanya merupakan upaya 
untuk menjelaskan konsep interval kepercayaan.  Apakah kedua pernyataan tersebut
mengatakan hal yang sama?  Jika tidak, uraian mana yang menurut Anda 
lebih akurat? 

\end{LJSItem}

\section[Percobaan Independen Diskret]{\bf Teorema Limit Pusat untuk Percobaan Independen Diskret}\label{sec 9.3}
Kita telah mengilustrasikan Teorema Limit Pusat untuk kasus percobaan Bernoulli, tetapi
teorema ini berlaku bagi kelas proses acak yang jauh lebih umum.  Secara khusus,
teorema ini berlaku bagi setiap proses percobaan independen yang setiap percobaannya
memiliki varians berhingga.  Untuk proses semacam itu, aproksimasi normal bagi suku-suku
individual maupun Teorema Limit Pusat berlaku.

Misalkan $S_n = X_1 + X_2 +\cdots+ X_n$ adalah jumlah $n$ peubah acak diskret
independen dari suatu proses percobaan independen, dengan fungsi distribusi yang sama $m(x)$ yang
didefinisikan pada bilangan bulat, rataan~$\mu$, dan varians~$\sigma^2$. \choice{Sebagaimana
telah kita lihat untuk jumlah-jumlah percobaan Bernoulli, distribusi untuk jumlah independen yang
lebih umum memiliki bentuk yang menyerupai kurva normal, tetapi nilai maksimumnya bergeser ke
kanan dan semakin mendatar.} {Dalam Bagian~\ref{sec 7.2}, kita telah melihat bahwa distribusi
untuk jumlah independen semacam itu memiliki bentuk yang menyerupai kurva normal, tetapi nilai
maksimumnya bergeser ke kanan dan kurvanya semakin mendatar (lihat Gambar~\ref{fig 7.9}).}
Kita dapat mencegah hal ini sebagaimana yang kita lakukan untuk percobaan Bernoulli.

\subsection*{Jumlah Terstandardisasi}
Perhatikan peubah acak terstandardisasi
$$
S_n^* = \frac {S_n - n\mu}{\sqrt{n\sigma^2}}\ .
$$

Standardisasi ini membuat $S_n$ memiliki nilai harapan~0 dan varians~1.  Jika $S_n
= j$, maka $S_n^*$ bernilai $x_j$, dengan
$$
x_j = \frac {j - n\mu}{\sqrt{n\sigma^2}}\ .
$$
Kita dapat membuat grafik paku sebagaimana yang kita lakukan untuk percobaan Bernoulli.  Setiap
paku berpusat di suatu~$x_j$.  Jarak antara dua paku berturutan adalah
$$
b = \frac 1{\sqrt{n\sigma^2}}\ ,
$$
dan tinggi paku tersebut adalah
$$
h = \sqrt{n\sigma^2} P(S_n = j)\ .
$$

Kasus percobaan Bernoulli merupakan kasus khusus dengan $X_j = 1$ jika hasil ke-$j$
merupakan keberhasilan dan 0 jika tidak; dalam hal ini $\mu = p$ dan $\sigma^2 = \sqrt {pq}$.
\par 
Sekarang kita mengilustrasikan proses ini untuk dua distribusi diskret yang berbeda.  Distribusi
pertama adalah $m$, yang diberikan oleh
$$
m = \pmatrix{
1 & 2 & 3 & 4 & 5 \cr
.2 & .2 & .2 & .2 & .2\cr}\ .
$$

Dalam Gambar~\ref{fig 9.5}, kita menampilkan jumlah terstandardisasi untuk distribusi ini dalam
kasus $n = 2$ dan $n = 10$.  Bahkan untuk $n = 2$, aproksimasinya ternyata sangat baik.

\putfig{5truein}{PSfig9-5}{Distribusi jumlah terstandardisasi.}{fig 9.5} 

\par
Untuk distribusi diskret kedua, kita memilih
$$
m = \pmatrix{
1 & 2 & 3 & 4 & 5 \cr
.4 & .3 & .1 & .1 & .1\cr}\ .
$$
\par
Distribusi ini cukup asimetris dan aproksimasinya tidak terlalu baik untuk $n
= 3$, tetapi pada $n = 10$ kita kembali memperoleh aproksimasi yang sangat baik (lihat
Gambar~\ref{fig 9.5.5}). Gambar~\ref{fig 9.5} dan \ref{fig 9.5.5} dibuat dengan program
{\bf CLTIndTrialsPlot}.
\index{CLTIndTrialsPlot (program)}

\putfig{5truein}{PSfig9-5-5}{Distribusi jumlah terstandardisasi.}{fig 9.5.5} 

\subsection*{Teorema Aproksimasi}
Seperti pada kasus percobaan Bernoulli, grafik-grafik ini menyarankan teorema aproksimasi
berikut untuk peluang-peluang individual.

\begin{theorem}\label{thm 9.3.5}
Misalkan $X_1$,~$X_2$, \dots,~$X_n$ merupakan proses percobaan independen dan misalkan $S_n =
X_1 + X_2 +\cdots+ X_n$.  Andaikan pembagi persekutuan terbesar dari selisih semua nilai yang
dapat diambil oleh $X_j$ adalah~1.  Misalkan $E(X_j) = \mu$ dan $V(X_j) = \sigma^2$.  Maka untuk
$n$ yang besar,
$$
P(S_n = j) \sim \frac {\phi(x_j)}{\sqrt{n\sigma^2}}\ ,
$$
dengan $x_j = (j - n\mu)/\sqrt{n\sigma^2}$, dan $\phi(x)$ adalah densitas normal baku.
\end{theorem}

Program {\bf CLTIndTrialsLocal}\index{CLTIndTrialsLocal (program)} menerapkan aproksimasi ini.
Ketika kita menjalankan program ini untuk 6 kali pelemparan sebuah dadu dan meminta peluang bahwa
jumlah hasil lemparannya sama dengan 21, kita memperoleh nilai sebenarnya .09285 dan nilai
aproksimasi normal .09537.  Jika kita menjalankan program ini untuk 24 kali pelemparan sebuah dadu
dan meminta peluang bahwa jumlah hasil lemparannya sama dengan 72, kita memperoleh nilai sebenarnya
.01724 dan nilai aproksimasi normal .01705.  Hasil-hasil ini menunjukkan bahwa aproksimasi normal
tersebut cukup baik.

\subsection*{Teorema Limit Pusat untuk Proses Percobaan Independen Diskret}
Teorema Limit Pusat untuk proses percobaan independen diskret adalah sebagai berikut.

\begin{theorem}{\bf (Teorema Limit Pusat)}\label{thm 9.3.6}\index{Teorema Limit Pusat!untuk
proses percobaan independen diskret} 
Misalkan $S_n = X_1 + X_2 +\cdots+ X_n$ adalah jumlah $n$ peubah acak diskret independen
dengan distribusi yang sama dan memiliki nilai harapan~$\mu$ serta varians~$\sigma^2$.  Maka, untuk
$a < b$,
$$
\lim_{n \to \infty} P\left( a < \frac {S_n - n\mu}{\sqrt{n\sigma^2}} < b\right)
= \frac 1{\sqrt{2\pi}} \int_a^b e^{-x^2/2}\, dx\ .
$$
\end{theorem}
\par
\choice{\footnote {Bukti Teorema \ref{thm 9.3.5}~dan Teorema~\ref{thm
9.3.6} diberikan dalam Bagian 10.3 pada buku Grinstead--Snell edisi lengkap.}}{Kita akan
memberikan bukti Teorema \ref{thm 9.3.5}~dan Teorema~\ref{thm 9.3.6} dalam
Bagian~\ref{sec 10.3}.}  Di sini kita membahas beberapa contoh.

\subsection*{Contoh}
\begin{example}
Sebuah dadu dilempar 420 kali.  Berapakah peluang bahwa jumlah hasil lemparannya
berada di antara 1400~dan~1550?

Jumlah tersebut adalah peubah acak
$$
S_{420} = X_1 + X_2 +\cdots+ X_{420}\ ,
$$
dengan setiap $X_j$ memiliki distribusi
$$
m_X = \pmatrix{
1 & 2 & 3 & 4 & 5 & 6 \cr
1/6 & 1/6 & 1/6 & 1/6 & 1/6 & 1/6 \cr}
$$
Kita telah melihat bahwa $\mu = E(X) = 7/2$ dan $\sigma^2 = V(X) = 35/12$.  Jadi,
$E(S_{420}) = 420 \cdot 7/2 = 1470$, $\sigma^2(S_{420}) = 420 \cdot 35/12 =
1225$, dan $\sigma(S_{420}) = 35$.  Oleh karena itu,
\begin{eqnarray*}
P(1400 \leq S_{420} \leq 1550) &\approx& P\left(\frac {1399.5 - 1470}{35} \leq
S_{420}^* \leq \frac {1550.5 - 1470}{35} \right) \\
     &=& P(-2.01 \leq S_{420}^* \leq 2.30) \\
&\approx&  \NA(-2.01, 2.30) = .9670\ . 
\end{eqnarray*}
Perhatikan bahwa program {\bf CLTIndTrialsGlobal} dapat digunakan untuk menghitung
peluang-peluang ini.
\end{example}

\begin{example}\label{exam 9.9}
Rata-rata nilai\index{rata-rata nilai} seorang mahasiswa adalah rata-rata nilainya dalam 30
mata kuliah.  Nilai-nilai tersebut didasarkan pada kemungkinan 100 poin dan dicatat sebagai
bilangan bulat.  Andaikan bahwa, dalam setiap mata kuliah, pengajar membuat galat penilaian sebesar
$k$ dengan peluang $|p/k|$, dengan $k = \pm1$,~$\pm2$, $\pm3$, $\pm4$,~$\pm5$.  Dengan demikian,
peluang tidak terjadi galat adalah $1 - (137/30)p$.  (Parameter~$p$ menyatakan tingkat
ketidakakuratan penilaian pengajar.)  Jadi, dalam setiap mata kuliah, terdapat dua nilai bagi
mahasiswa tersebut, yaitu nilai yang ``benar" dan nilai yang tercatat.  Karena itu, mahasiswa
tersebut memiliki dua rata-rata nilai, yaitu rata-rata nilai yang benar dan rata-rata nilai yang
tercatat.
\par
Kita ingin memperkirakan peluang bahwa kedua rata-rata nilai ini berbeda kurang dari .05 untuk
seorang mahasiswa tertentu.  Sekarang kita andaikan bahwa $p = 1/20$.  Kita juga mengandaikan bahwa
galat total merupakan jumlah $S_{30}$ dari 30 peubah acak independen yang masing-masing berdistribusi
$$
m_X: \left\{ 
\begin{array}{ccccccccccc}
-5 & -4 & -3 & -2 & -1 & 0 & 1 & 2 & 3 & 4 & 5 \\
\frac1{100} & \frac1{80} & \frac1{60} & \frac1{40} & \frac1{20} & 
\frac{463}{600} & \frac1{20} & \frac1{40} & \frac1{60} & \frac1{80} &
\frac1{100}
\end{array} 
\right \}\ .
$$
Dapat dihitung dengan mudah bahwa $E(X) = 0$ dan $\sigma^2(X) = 1.5$.  Dengan demikian,
$$\begin{array}{ll}
P\left(-.05 \leq \frac {S_{30}}{30} \leq .05 \right) &= P(-1.5 \leq S_{30} \leq1.5) \\
     & \\
     &= P\left(\frac {-1.5}{\sqrt{30\cdot1.5}} \leq S_{30}^* \leq \frac {1.5}{\sqrt{30\cdot1.5}} \right) \\
     & \\
     &= P(-.224 \leq S_{30}^* \leq .224) \\
     & \\
     &\approx   \NA(-.224, .224) = .1772\ .  
\end{array}$$
Ini berarti bahwa peluang rata-rata nilai seorang mahasiswa tertentu akurat hingga .05
hanya sebesar 17.7\%.  (Jadi, misalnya, jika dua calon lulusan terbaik memiliki rata-rata tercatat
97.1 dan 97.2, terdapat peluang yang tidak dapat diabaikan bahwa urutan rata-rata mereka yang benar
justru terbalik.)  Untuk pembahasan lebih lanjut mengenai contoh ini, lihat artikel karya R.~M.
Kozelka.\index{KOZELKA, R. M.}\footnote{R. M. Kozelka, ``Grade-Point Averages and the Central Limit
Theorem," \emx {American Math.\ Monthly,} vol.~86 (Nov 1979), pp.~773-777.}
\end{example}



\subsection*{Teorema Limit Pusat yang Lebih Umum}
Dalam Teorema~\ref{thm 9.3.6}, peubah-peubah acak diskret yang dijumlahkan diasumsikan independen
dan berdistribusi identik.  Ternyata, asumsi distribusi identik dapat diperlemah secara berarti.
Banyak penelitian telah dilakukan dalam bidang ini, dengan salah satu sumbangan penting diberikan
oleh J. W. Lindeberg\index{LINDEBERG, J. W.}.  Lindeberg menemukan suatu syarat bagi barisan
$\{X_n\}$ yang menjamin bahwa distribusi jumlah $S_n$ berdistribusi normal secara asimtotik.
Feller menunjukkan bahwa syarat Lindeberg juga diperlukan, dalam arti bahwa jika syarat tersebut
tidak dipenuhi, maka jumlah $S_n$ tidak berdistribusi normal secara asimtotik.  Untuk pernyataan
Teorema Lindeberg yang tepat, pembaca kami rujuk kepada Feller.\index{FELLER, W.}\footnote{W.~Feller,
\emx {Introduction to Probability Theory and its Applications,} vol.~1, 3rd~ed. (New York: John
Wiley \& Sons, 1968), p.~254.}  Syarat cukup yang lebih kuat (tetapi lebih mudah dinyatakan)
daripada syarat Lindeberg, dan lebih lemah daripada syarat dalam Teorema~\ref{thm 9.3.6}, diberikan
dalam teorema berikut.
\pagebreak[4]
\begin{theorem}{\bf (Teorema Limit Pusat)}\label{thm 9.3.7}\index{Teorema Limit Pusat!untuk
peubah diskret independen} 
Misalkan $X_1,\ X_2,\ \ldots,\ X_n\ ,\ \ldots$ adalah barisan peubah acak diskret independen,
dan misalkan $S_n = X_1 + X_2 +\cdots+ X_n$.  Untuk setiap $n$, nyatakan rataan dan varians $X_n$
berturut-turut dengan $\mu_n$ dan $\sigma^2_n$.  Definisikan rataan dan varians $S_n$
berturut-turut sebagai $m_n$ dan $s_n^2$, dan andaikan bahwa $s_n \rightarrow \infty$.  Jika
terdapat suatu konstanta $A$ sedemikian sehingga $|X_n| \le A$ untuk semua $n$, maka untuk $a < b$,
$$
\lim_{n \to \infty} P\left( a < \frac {S_n - m_n}{s_n} < b\right)
= \frac 1{\sqrt{2\pi}} \int_a^b e^{-x^2/2}\, dx\ .
$$
\end{theorem}
Syarat bahwa $|X_n| \le A$ untuk semua $n$ kadang-kadang dinyatakan dengan mengatakan bahwa
barisan $\{X_n\}$ terbatas seragam.  Syarat bahwa $s_n \rightarrow \infty$ diperlukan (lihat
Latihan~\ref{exer 9.2.114}).
\par
Kita mengilustrasikan teorema ini dengan membangkitkan barisan $n$ distribusi acak pada
interval $[a, b]$.  Kemudian kita mengonvolusikan distribusi-distribusi ini untuk mencari distribusi
jumlah dari $n$ percobaan independen yang diatur oleh distribusi tersebut.  Terakhir, kita
menstandardisasi distribusi jumlah itu agar memiliki rataan~0 dan simpangan baku~1, lalu
membandingkannya dengan densitas normal.  Program {\bf CLTGeneral}\index{CLTGeneral (program)}
menjalankan prosedur ini.

\putfig{3.5truein}{PSfig9-6}{Jumlah peubah acak yang dipilih secara acak.}{fig 9.6} 

Dalam Gambar~\ref{fig 9.6}, kita menampilkan hasil menjalankan program ini untuk
$[a, b] = [-2, 4]$, dan $n = 1,\ 4,$ dan 10.  Terlihat bahwa distribusi acak pertama kita cukup
asimetris.  Ketika kita memilih jumlah sepuluh percobaan semacam itu, kecocokannya dengan kurva
normal sudah sangat baik.
\par
Pada dasarnya, teorema di atas mengatakan bahwa apa pun yang dapat dipandang tersusun sebagai
jumlah banyak bagian kecil yang independen berdistribusi mendekati normal.  Hal ini membawa kita
kepada salah satu pertanyaan terpenting yang diajukan mengenai genetika pada tahun 1800-an.


\subsection*{Distribusi Normal dan Genetika}\index{genetika}
Jika kita mengamati distribusi tinggi badan\index{tinggi badan!distribusi} orang dewasa dari satu
jenis kelamin dalam suatu populasi, sulit untuk tidak memperhatikan bahwa distribusi ini tampak
seperti distribusi normal.  Contohnya ditampilkan dalam Gambar~\ref{fig 9.61}.  Gambar ini
menampilkan distribusi tinggi badan 9593 perempuan berusia antara 21 dan 74 tahun.  Data ini
berasal dari Health and Nutrition Examination Survey I (HANES I).\index{data HANES}
Untuk survei ini, dipilih suatu sampel dari penduduk sipil Amerika Serikat.  Survei tersebut
dilaksanakan antara tahun 1971 dan 1974.
\par
Pertanyaan yang wajar diajukan ialah ``Bagaimana hal ini dapat terjadi?".  Francis
Galton,\index{GALTON, F.} ilmuwan Inggris pada abad ke-19, mempelajari pertanyaan ini dan berbagai
pertanyaan terkait, serta menyusun model-model peluang yang sangat penting untuk menjelaskan
pengaruh genetis pada atribut seperti tinggi badan.  Bahkan, salah satu gagasan terpenting dalam
statistika, yaitu regresi menuju rataan\index{regresi menuju rataan}, ditemukan oleh Galton dalam
upayanya memahami pengaruh genetis tersebut.
 
\putfig{3.5truein}{PSfig9-61}{Distribusi tinggi badan perempuan dewasa.}{fig 9.61} 
\par
Galton menghadapi sebuah kontradiksi yang tampak nyata.  Di satu sisi, ia mengetahui bahwa
distribusi normal muncul dalam situasi ketika banyak pengaruh kecil yang independen dijumlahkan.
Di sisi lain, ia juga mengetahui bahwa banyak atribut kuantitatif, seperti tinggi badan, sangat
dipengaruhi oleh faktor genetis: orang tua yang tinggi cenderung memiliki keturunan yang tinggi.
Jadi, dalam kasus ini tampaknya terdapat dua pengaruh besar, yaitu kedua orang tua.  Galton tentu
menyadari bahwa faktor non-genetis ikut berperan dalam menentukan tinggi badan seseorang.  Namun,
kecuali jika faktor non-genetis ini mengalahkan faktor genetis---yang berarti membantah hipotesis
bahwa keturunan penting dalam menentukan tinggi badan---tampaknya mustahil bagi pasangan orang tua
dengan tinggi badan tertentu untuk memiliki keturunan yang tinggi badannya berdistribusi normal.
\par
Masalah di atas dapat dinyatakan secara simbolis sebagai berikut.  Andaikan kita memilih dua
bilangan riil positif tertentu $x$ dan $y$, lalu mencari semua pasangan orang tua dengan salah
satunya memiliki tinggi $x$ satuan dan yang lainnya $y$ satuan.  Kemudian kita mengamati seluruh
keturunan pasangan-pasangan tersebut.  Kita dapat mempostulatkan adanya fungsi $f(x, y)$ yang
menyatakan pengaruh genetis tinggi badan orang tua terhadap tinggi badan keturunan.  Selanjutnya,
kita dapat membiarkan $W$ menyatakan pengaruh faktor non-genetis terhadap tinggi badan keturunan.
Maka, untuk pasangan tinggi badan tertentu $\{x, y\}$, peubah acak yang menyatakan tinggi badan
keturunan diberikan oleh
$$H = f(x, y) + W\ ,$$
dengan $f$ merupakan fungsi deterministik, yaitu fungsi yang memberikan satu keluaran untuk
pasangan masukan $\{x, y\}$.  Jika kita mengandaikan bahwa pengaruh $f$ besar dibandingkan dengan
pengaruh $W$, maka varians $W$ kecil.  Namun, karena f deterministik, varians $H$ sama dengan
varians $W$, sehingga varians $H$ kecil.  Akan tetapi, dari datanya Galton mengamati bahwa varians
tinggi badan keturunan dari pasangan tinggi badan orang tua tertentu tidaklah kecil.  Hal ini
tampaknya menyiratkan bahwa pewarisan hanya berperan kecil dalam menentukan tinggi badan seseorang.
Nanti dalam seksi ini, kita akan menguraikan cara Galton mengatasi masalah tersebut.
\par
Sekarang kita akan membahas penjelasan modern mengenai mengapa sifat-sifat tertentu, seperti tinggi
badan, berdistribusi mendekati normal.  Untuk itu, kita perlu memperkenalkan beberapa istilah dari
bidang genetika.  Sel-sel\index{sel} dalam organisme hidup yang tidak terlibat langsung dalam
pewarisan materi genetis kepada keturunan disebut sel somatik, sedangkan sel-sel lainnya disebut
sel germinal.  Informasi genetis organisme dari suatu spesies dikodekan dalam kumpulan entitas
fisik yang disebut kromosom\index{kromosom}.  Kromosom-kromosom berpasangan dalam setiap sel
somatik.  Sebagai contoh, manusia memiliki 23 pasang kromosom dalam setiap sel somatik.  Sel
kelamin memuat satu kromosom dari setiap pasangan.  Dalam reproduksi seksual, dua sel kelamin,
masing-masing dari satu orang tua, menyumbangkan kromosomnya untuk membentuk kumpulan kromosom
keturunan.
\par
Kromosom memuat banyak subunit yang disebut gen\index{gen}.  Gen tersusun dari molekul
DNA\index{DNA}, dan DNA suatu gen mengodekan informasi yang mengatur protein.  Dalam konteks ini,
kita akan membahas gen-gen yang memuat informasi yang memengaruhi suatu sifat fisik organisme,
seperti tinggi badan.  Kromosom yang berpasangan membuat gen-gen pada kromosom itu turut berpasangan.
\par
Dalam suatu spesies, setiap gen dapat memiliki salah satu dari beberapa bentuk.  Berbagai bentuk
ini disebut alel\index{alel}.  Alel-alel yang berbeda dapat dipandang berpotensi menimbulkan
pengaruh yang berbeda pada sifat fisik yang dibahas.  Dari dua alel yang terdapat dalam suatu
pasangan gen pada sebuah organisme, satu alel berasal dari salah satu orang tua dan alel lainnya
berasal dari orang tua yang lain.  Jenis-jenis pasangan alel yang mungkin (tanpa memperhatikan
urutan) disebut genotipe\index{genotipe}.
\par
Jika kita mengandaikan bahwa tinggi badan manusia sebagian besar dikendalikan oleh sebuah gen
tertentu, kita menghadapi kesulitan yang sama seperti Galton.  Kita mengandaikan bahwa setiap orang
tua memiliki sepasang alel yang sebagian besar mengendalikan tinggi badannya.  Karena setiap orang
tua menyumbangkan satu alel dari pasangan gen ini kepada setiap keturunannya, terdapat empat
pasangan alel yang mungkin bagi keturunan pada lokasi gen ini.  Asumsinya ialah bahwa pasangan alel
tersebut sebagian besar mengendalikan tinggi badan keturunan, dan kita juga mengandaikan bahwa
faktor genetis lebih berpengaruh daripada faktor non-genetis.  Akibatnya, pada keturunan seharusnya
terlihat beberapa modus dalam distribusi tinggi badan, satu modus untuk setiap pasangan alel yang
mungkin.  Distribusi ini tidak sesuai dengan distribusi tinggi badan yang diamati.
\par
Hipotesis alternatif yang dapat menjelaskan pengamatan bahwa tinggi badan keturunan dengan jenis
kelamin tertentu berdistribusi normal ialah hipotesis banyak-gen\index{hipotesis banyak-gen}.
Menurut hipotesis ini, kita mengandaikan bahwa banyak gen memengaruhi tinggi badan seseorang.
Besarnya pengaruh gen-gen tersebut dapat berbeda-beda.  Karena itu, kita dapat merepresentasikan
setiap pasangan gen dengan peubah acak $X_i$, dengan nilai peubah acak tersebut merupakan pengaruh
pasangan alel terhadap tinggi badan orang itu.  Sebagai contoh, jika setiap orang tua memiliki dua
alel berbeda dalam pasangan gen yang sedang dibahas, maka keturunannya memiliki salah satu dari
empat pasangan alel yang mungkin pada lokasi gen ini.  Sekarang tinggi badan keturunan merupakan
peubah acak yang dapat dinyatakan sebagai
$$H = X_1 + X_2 + \cdots + X_n + W\ ,$$
jika terdapat $n$ gen yang memengaruhi tinggi badan.  (Di sini, seperti sebelumnya, peubah acak
$W$ menyatakan pengaruh non-genetis.)  Meskipun $n$ tetap, jika nilainya cukup besar, maka
Teorema~\ref{thm 9.3.7} menyiratkan bahwa jumlah $X_1 + X_2 + \cdots + X_n$ berdistribusi mendekati
normal.  Selanjutnya, jika kita mengandaikan bahwa pengaruh kumulatif $X_i$ jauh lebih besar
daripada pengaruh $W$, maka $H$ berdistribusi mendekati normal.
\par
Ciri lain yang diamati dari distribusi tinggi badan orang dewasa dengan satu jenis kelamin dalam
suatu populasi ialah bahwa variansnya tampak tidak bertambah atau berkurang dari satu generasi ke
generasi berikutnya.  Hal ini telah diketahui pada masa Galton, dan upayanya untuk menjelaskannya
membawanya kepada gagasan regresi menuju rataan.  Gagasan ini akan dibahas lebih lanjut dalam
catatan sejarah di akhir seksi.  (Alasan kita hanya membahas satu jenis kelamin ialah bahwa tinggi
badan manusia jelas terkait dengan jenis kelamin; secara umum, jika kita memiliki dua populasi
yang masing-masing berdistribusi normal, gabungan keduanya belum tentu berdistribusi normal.)
\par
Dengan hipotesis banyak-gen, mudah untuk menjelaskan mengapa varians seharusnya tetap dari satu
generasi ke generasi berikutnya.  Kita mulai dengan mengandaikan bahwa pada suatu lokasi gen
tertentu terdapat $k$ alel, yang akan kita nyatakan dengan $A_1,\ A_2,\ \ldots,\ A_k$.  Kita
mengandaikan bahwa keturunan dihasilkan melalui perkawinan acak.  Maksudnya, setiap keturunan
memiliki peluang yang sama untuk berasal dari pasangan orang tua mana pun dalam generasi
sebelumnya.  Ada cara lain memandang perkawinan acak yang mempermudah perhitungan.  Perhatikan
himpunan $S$ yang memuat seluruh alel (pada lokasi gen yang ditentukan) dalam semua sel germinal
seluruh individu pada generasi orang tua.  Dalam konteks himpunan $S$, perkawinan acak berarti
setiap pasangan alel dalam $S$ berpeluang sama untuk berada dalam keturunan tertentu.  (Pembaca
mungkin berkeberatan terhadap cara memandang perkawinan acak ini karena memungkinkan dua alel dari
orang tua yang sama berakhir pada satu keturunan; tetapi jika banyak individu dalam populasi orang
tua cukup besar, diizinkan atau tidaknya peristiwa ini tidak banyak memengaruhi peluang.)
\par
Untuk $1 \le i \le k$, misalkan $p_i$ menyatakan proporsi alel dalam populasi orang tua yang
bertipe $A_i$.  Jelas bahwa proporsi ini sama dengan proporsi alel dalam sel germinal populasi
orang tua, dengan mengandaikan bahwa setiap orang tua menghasilkan sel germinal dalam jumlah yang
kurang lebih sama.  Perhatikan distribusi alel pada keturunan.  Karena setiap sel germinal berpeluang sama
untuk dipilih bagi keturunan tertentu, distribusi alel pada keturunan sama dengan distribusi pada
orang tua.
\par
Selanjutnya, kita membahas distribusi genotipe dalam kedua generasi.  Kita akan membuktikan fakta
berikut: distribusi genotipe dalam generasi keturunan hanya bergantung pada distribusi alel dalam
generasi orang tua (khususnya, tidak bergantung pada distribusi genotipe dalam generasi orang tua).
Perhatikan genotipe-genotipe yang mungkin; terdapat $k(k+1)/2$ genotipe.  Berdasarkan asumsi kita,
genotipe $A_iA_i$ akan muncul dengan frekuensi $p_i^2$, dan genotipe $A_iA_j$, dengan $i \ne j$,
akan muncul dengan frekuensi $2p_ip_j$.  Jadi, seperti yang dinyatakan, frekuensi genotipe hanya
bergantung pada frekuensi alel dalam generasi orang tua.
\par
Ini berarti bahwa jika kita mulai dari generasi tertentu dan distribusi alel tertentu, maka dalam
semua generasi sesudah generasi awal tersebut, baik distribusi alel maupun distribusi genotipe akan
tetap.  Pernyataan terakhir ini dikenal sebagai Hukum Hardy--Weinberg.\index{Hukum Hardy--Weinberg}
\par
Kita dapat menguraikan akibat hukum ini bagi distribusi tinggi badan orang dewasa dari satu jenis
kelamin dalam suatu populasi.  Ingat bahwa tinggi badan keturunan diberikan oleh peubah acak $H$,
dengan
$$H = X_1 + X_2 + \cdots + X_n + W\ ,$$
dengan $X_i$ bersesuaian dengan gen-gen yang memengaruhi tinggi badan, dan peubah acak $W$
menyatakan pengaruh non-genetis.  Hukum Hardy--Weinberg menyatakan bahwa untuk setiap $X_i$,
distribusi pada generasi keturunan sama dengan distribusi pada generasi orang tua.  Jadi, jika kita
mengandaikan bahwa distribusi $W$ kurang lebih sama dari satu generasi ke generasi berikutnya (atau
jika kita mengandaikan bahwa pengaruhnya kecil), maka distribusi $H$ sama dari satu generasi ke
generasi berikutnya.  (Sesungguhnya, pengaruh pola makan merupakan bagian dari $W$, dan jelas bahwa
dalam banyak populasi manusia, pola makan telah cukup banyak berubah antargenerasi pada masa
belakangan ini.  Perubahan ini dianggap sebagai salah satu alasan mengapa manusia rata-rata semakin
tinggi.  Pengaruh $W$ juga dianggap kecil dibandingkan dengan pengaruh genetis orang tua.)

\subsection*{Pembahasan}
Secara umum, Teorema Limit Pusat memuat lebih banyak informasi daripada Hukum Bilangan Besar,
karena teorema tersebut memberikan informasi terperinci mengenai \emx {bentuk} distribusi $S_n^*$;
untuk~$n$ yang besar, bentuknya kurang lebih sama dengan bentuk densitas normal baku.  Secara lebih
khusus, Teorema Limit Pusat menyatakan bahwa jika kita menstandardisasi distribusi $S_n$ dan
mengoreksi tingginya, maka fungsi densitas normal merupakan aproksimasi yang sangat baik bagi
distribusi ini ketika $n$ besar.  Dengan demikian, kita memiliki aproksimasi yang dapat dihitung
untuk distribusi~$S_n$.  Aproksimasi ini memberi kita teknik yang ampuh untuk menghasilkan jawaban
atas berbagai macam pertanyaan mengenai jumlah peubah acak independen, sekalipun setiap peubah acak
memiliki distribusi yang berbeda.

\subsection*{Catatan Sejarah}
Pada pertengahan tahun 1800-an, matematikawan Belgia Quetelet\index{QUETELET,
A.}\footnote{S. Stigler, \emx {The History of Statistics,} (Cambridge: Harvard University Press,
1986), p. 203.} telah menunjukkan secara empiris bahwa distribusi normal muncul dalam data nyata,
dan juga telah memberikan metode untuk mencocokkan kurva normal dengan suatu kumpulan data.
Jauh sebelumnya, Laplace\index{LAPLACE, P. S.}\footnote{ibid., p. 136} telah menunjukkan bahwa
jumlah banyak peubah acak independen yang berdistribusi identik mendekati normal.
Galton\index{GALTON, F.} mengetahui bahwa sifat fisik tertentu dalam suatu populasi tampak
berdistribusi mendekati normal, tetapi ia tidak menganggap hasil Laplace sebagai penjelasan yang
baik mengenai bagaimana distribusi ini muncul.  Berikut ini kita kutip pernyataan Galton dari
sumber primer ranah publik.\footnote{Francis Galton, ``Typical Laws of Heredity,''
\emx{Nature}, vol.~15 (1877), pp.~512--514, doi:10.1038/015512b0. Terjemahan editorial
atas sumber primer ranah publik.}
\begin{quote}
Pertama-tama, izinkan saya menunjukkan suatu fakta yang entah mengapa luput dari perhatian
Quetelet dan semua penulis yang mengikuti jejaknya, dan yang berkaitan erat dengan pekerjaan kita
malam ini.  Faktanya ialah bahwa, meskipun ciri-ciri tumbuhan dan hewan mengikuti hukum tersebut,
alasan mengapa demikian masih sama sekali belum dijelaskan.  Inti hukum itu adalah bahwa
perbedaan-perbedaan harus sepenuhnya disebabkan oleh tindakan bersama sejumlah besar pengaruh
\emx{kecil} yang independen dalam berbagai kombinasi...Namun proses pewarisan...bukanlah pengaruh
kecil, melainkan pengaruh yang sangat penting...Kesimpulannya ialah...bahwa proses pewarisan harus
bekerja selaras dengan hukum penyimpangan, dan dalam pengertian tertentu juga harus mematuhinya.
\end{quote}
\par
Galton menciptakan alat yang dikenal sebagai quincunx\index{quincunx} (sekarang lazim disebut
papan Galton)\index{papan Galton}, yang kita gunakan dalam Contoh~\ref{exam 3.2.1} untuk
menunjukkan cara memperoleh distribusi binomial secara fisik.  Tentu saja, Teorema Limit Pusat
menyatakan bahwa untuk nilai parameter $n$ yang besar, distribusi binomial mendekati normal.
Galton menggunakan quincunx untuk menjelaskan bagaimana pewarisan memengaruhi distribusi suatu
sifat pada keturunan.
\par
Seperti Galton, kita membahas apa yang terjadi jika, pada suatu ketinggian di tengah, kita menghentikan
sementara gerak butiran yang jatuh di dalam quincunx.  Pembaca dapat melihat
skema konseptual pada Gambar~\ref{fig 9.62}.
\begin{figure}[ht]
\centering
\fbox{\begin{minipage}{4truein}
\centering
\small
Distribusi butiran pada garis $AB$: $N(\mu,\sigma_1^2)$\\[.15in]
$\downarrow$ setiap kompartemen dilepas dan butirannya kembali menyebar\\[.15in]
Distribusi jatuh dari satu kompartemen: $N(0,\sigma_2^2)$\\[.15in]
$\downarrow$ campuran semua kompartemen\\[.15in]
Distribusi akhir: $N(\mu,\sigma_1^2+\sigma_2^2)$
\end{minipage}}
\caption{Skema konseptual quincunx dua tahap, digambar ulang secara mandiri.}
\label{fig 9.62}
\end{figure}
Skema ini dirumuskan ulang berdasarkan uraian Karl Pearson,\index{PEARSON,
K.}\footnote{Rujukan historis: Karl Pearson, \emx {The Life, Letters and
Labours of Francis Galton,} vol. IIIB (Cambridge at the University Press,
1930), p.~466. Gambar historis tidak direproduksi.} yang bertumpu pada catatan
Galton.  Dalam skema ini, butiran untuk sementara
dipisahkan ke dalam kompartemen-kompartemen pada garis AB.  (Garis
A$^{\prime}$B$^{\prime}$ membentuk landasan tempat butiran dapat berhenti.)  Jika garis AB tidak
terlalu dekat dengan bagian atas quincunx, maka butiran akan berdistribusi mendekati normal pada
garis ini.  Sekarang andaikan satu kompartemen dibuka, seperti ditunjukkan dalam gambar.  Butiran
dari kompartemen tersebut akan jatuh dan membentuk distribusi normal di bagian bawah quincunx.
Jika kemudian semua kompartemen dibuka, seluruh butiran akan jatuh dan menghasilkan distribusi
yang sama seperti jika butiran itu tidak dihentikan sementara pada garis AB.  Namun, tindakan
menghentikan butiran pada garis AB lalu melepaskan kompartemen satu per satu sama saja dengan
mengonvolusikan dua distribusi normal.  Distribusi-distribusi normal di bagian bawah, yang
bersesuaian dengan setiap kompartemen pada garis AB, sedang dicampur, dengan bobot berupa banyaknya
butiran dalam setiap kompartemen.  Di sisi lain, telah diketahui bahwa jika butiran tidak
terhalang, distribusi akhirnya mendekati normal.  Jadi, alat ini menunjukkan bahwa konvolusi dua
distribusi normal kembali merupakan distribusi normal.
\par
Galton juga meninjau quincunx dari sudut pandang lain.  Ia memilah sekumpulan 490 biji ercis manis
berdasarkan beratnya ke dalam tujuh kelompok.  Kepada masing-masing dari tujuh temannya, ia
memberikan 10 biji dari setiap kelompok; teman-temannya kemudian menanam biji-biji itu.  Galton
menemukan bahwa setiap kelompok menghasilkan biji yang beratnya berdistribusi normal.  (Ercis
manis bereproduksi melalui penyerbukan sendiri, sehingga ia tidak perlu mempertimbangkan
kemungkinan interaksi antara kelompok yang berbeda.)  Selain itu, ia menemukan bahwa varians berat
keturunan sama untuk setiap kelompok.  Pemilahan ke dalam kelompok ini bersesuaian dengan
kompartemen pada garis AB dalam quincunx.  Jadi, ercis manis tersebut seolah-olah diatur oleh
konvolusi distribusi normal.
\par
Kini ia menghadapi suatu masalah.  Kita telah menunjukkan dalam Bab~\ref{chp 7}, dan Galton pun
mengetahuinya, bahwa konvolusi dua distribusi normal menghasilkan distribusi normal dengan varians
yang lebih besar daripada salah satu distribusi asal.  Namun, data biji ercis manisnya menunjukkan
bahwa varians populasi keturunan sama dengan varians populasi orang tua.  Jawabannya terhadap
masalah ini ialah mempostulatkan mekanisme yang ia sebut \emx{reversi}\index{reversi}, dan yang
sekarang disebut \emx{regresi menuju rataan}\index{regresi menuju rataan}.  Menurut ringkasan
Stigler,\index{STIGLER, S.}\footnote{S. Stigler, \emx{The History of Statistics}
(Cambridge: Harvard University Press, 1986), p.~282. Ringkasan editorial; ungkapan sumber
tidak diterjemahkan atau direproduksi.} rataan setiap kelompok keturunan bergeser menuju
rataan populasi; simpangannya tetap searah dengan simpangan kelompok induk, tetapi hanya sekitar
sepertiganya. Penyusutan linear ini, bersama variasi dalam setiap kelompok, mempertahankan
keragaman populasi.
\par
Mekanisme reversi Galton diringkas oleh skema konseptual pada
Gambar~\ref{fig 9.63}.\footnote{Rujukan historis: Karl Pearson, \emx {The Life,
Letters and Labours of Francis Galton,} vol. IIIA (Cambridge at the University
Press, 1930), p.~9. Gambar historis tidak direproduksi.}
\begin{figure}[ht]
\centering
\fbox{\begin{minipage}{4truein}
\centering
\small
Rataan kelompok orang tua: $\mu+d$\\[.15in]
$\downarrow\quad$ reversi linear $d\longmapsto d/3$\\[.15in]
Rataan kelompok keturunan: $\mu+d/3$\\[.15in]
Variasi di sekitar rataan kelompok menghasilkan kembali sebaran populasi.
\end{minipage}}
\caption{Skema mandiri penjelasan Galton mengenai reversi menuju rataan.}
\label{fig 9.63}
\end{figure}

%Should we put some exercises about genetics or Hardy-Weinberg?

\exercises
\begin{LJSItem}

\i\label{exer 9.2.100}  Sebuah dadu dilempar 24 kali.  Gunakan Teorema Limit Pusat untuk
memperkirakan peluang bahwa
\begin{enumerate}
\item  jumlah hasil lemparannya lebih besar dari 84.

\item  jumlah hasil lemparannya sama dengan 84.
\end{enumerate}

\i\label{exer 9.2.101}  Seorang pejalan acak mulai dari~0 pada sumbu~$x$ dan pada setiap satuan
waktu bergerak 1 langkah ke kanan atau 1 langkah ke kiri dengan peluang 1/2.  Perkirakan peluang
bahwa, setelah 100 langkah, pejalan tersebut berada lebih dari 10 langkah dari posisi awal.

\i\label{exer 9.2.102}  Seutas tali tersusun dari 100 serat.  Andaikan bahwa kekuatan putus tali
tersebut merupakan jumlah kekuatan putus setiap serat.  Andaikan pula bahwa jumlah ini dapat
dipandang sebagai jumlah suatu proses percobaan independen dengan 100 percobaan yang masing-masing
memiliki nilai harapan 10 pon dan simpangan baku~1.  Tentukan peluang hampiran bahwa tali tersebut
dapat menahan beban
\begin{enumerate}
\item  1000 pon.

\item  970 pon.
\end{enumerate}

\i\label{exer 9.2.103}  Tulislah sebuah program untuk mencari rata-rata 1000 digit acak
0,~1, 2, 3, 4, 5, 6, 7, 8, atau~9.  Buatlah program tersebut menguji apakah rata-ratanya berada
dalam jarak tiga simpangan baku dari nilai harapan~4.5.  Ubahlah program agar mengulangi simulasi
ini 1000 kali dan mencatat berapa kali pengujian berhasil.  Apakah hasil Anda sesuai dengan
Teorema Limit Pusat?

\i\label{exer 9.2.104}  Sebuah dadu dilempar hingga untuk pertama kalinya jumlah seluruh nilai sisi
yang muncul mencapai 700 atau lebih.  Perkirakan peluang bahwa untuk mencapai jumlah tersebut diperlukan
\begin{enumerate}
\item  lebih dari 210 lemparan.

\item  kurang dari 190 lemparan.

\item  antara 180 dan 210 lemparan, termasuk kedua batasnya.
\end{enumerate}

\i\label{exer 9.2.105}  Sebuah bank menerima gulungan koin satu sen dan mengkreditkan 50~sen kepada
nasabah tanpa menghitung isinya.  Andaikan bahwa sebuah gulungan berisi 49 koin dalam
30~persen kasus, 50 koin dalam 60~persen kasus, dan 51 koin dalam 10~persen kasus.
\begin{enumerate}
\item  Tentukan nilai harapan dan varians besar kerugian bank pada sebuah gulungan biasa.

\item  Perkirakan peluang bahwa bank akan kehilangan lebih dari 25~sen dalam 100 gulungan.

\item  Perkirakan peluang bahwa bank akan kehilangan tepat 25~sen dalam 100 gulungan.

\item  Perkirakan peluang bahwa bank akan mengalami kerugian berapa pun dalam 100 gulungan.

\item  Berapa banyak gulungan yang perlu dikumpulkan bank agar bank memiliki peluang 99 persen
mengalami kerugian bersih?
\end{enumerate}

\i\label{exer 9.2.106}  Ketika mengukur tinggi sebuah menara setinggi 200 kaki, suatu alat ukur
survei menghasilkan galat $-2$,~$-1$, 0, 1, atau~2 kaki dengan peluang yang sama.
\begin{enumerate}
\item  Tentukan nilai harapan dan varians tinggi yang diperoleh dari satu kali penggunaan alat ini.

\item  Perkirakan peluang bahwa dalam 18 pengukuran independen terhadap menara ini, rata-rata hasil
pengukuran berada di antara 199~dan~201, termasuk kedua batasnya.
\end{enumerate}

\i\label{exer 9.2.107}  Untuk Contoh~\ref{exam 9.9}, perkirakan $P(S_{30} = 0)$.  Artinya,
perkirakan peluang bahwa galat-galatnya saling meniadakan dan rata-rata nilai mahasiswa tersebut
benar.

\i\label{exer 9.2.108}  Buktikan Hukum Bilangan Besar dengan menggunakan Teorema Limit Pusat.

\i\label{exer 9.2.109}  Peter dan Paul memainkan permainan cocok koin 10{,}000 kali.  Uraikan
secara singkat apa yang dinyatakan oleh setiap teorema berikut mengenai kekayaan Peter.
\begin{enumerate}
\item  Hukum Bilangan Besar.

\item  Teorema Limit Pusat.
\end{enumerate}

\i\label{exer 9.2.110}  Seorang wisatawan di Las Vegas tertarik pada suatu permainan judi yang
meminta pemain mempertaruhkan 1~dolar pada setiap permainan; kemenangan membayar pemain 2~dolar
ditambah pengembalian taruhannya, sedangkan kekalahan hanya membuatnya kehilangan taruhan.
Orang dalam Las Vegas dan pelajar teori peluang yang jeli mengetahui bahwa peluang menang dalam
permainan ini adalah~1/4.  Ketika rasa lapar memaksanya meninggalkan meja, wisatawan tersebut telah
memainkan permainan ini 240 kali.  Dengan mengandaikan bahwa tidak terjadi sesuatu yang nyaris
mustahil, kira-kira berapa banyak uang yang hilang ketika wisatawan itu meninggalkan kasino?
Berapakah peluang bahwa ia tidak kehilangan uang?

\i\label{exer 9.2.111}  Kita telah melihat bahwa, ketika bermain rolet di Monte Carlo
(Contoh~\ref {exam 6.7}), bertaruh 1~dolar pada merah atau 1~dolar pada~17 berarti memilih di antara
distribusi-distribusi
$$
m_X = \pmatrix{
-1 & -1/2 & 1 \cr
18/37 & 1/37 & 18/37\cr }
$$
atau
$$
m_X = \pmatrix{
-1 & 35 \cr
36/37 & 1/37 \cr }
$$
Anda berencana memilih salah satu metode ini dan menggunakannya untuk memasang 100 taruhan
masing-masing sebesar 1 dolar.  Dengan menggunakan Teorema Limit Pusat, perkirakan peluang memperoleh
keuntungan dalam masing-masing permainan.  Bandingkan perkiraan Anda dengan peluang sebenarnya,
yang melalui perhitungan eksak dapat ditunjukkan bernilai .437 dan .509 hingga tiga tempat desimal.

\i\label{exer 9.2.112}  Dalam Contoh~\ref{exam 9.9}, tentukan nilai terbesar~$p$ yang memberikan
peluang .954 bahwa angka pada tempat desimal pertama benar.

\i\label{exer 9.2.113}  Contoh~\ref{exam 9.9} dianggap tidak realistis karena peluang galatnya
terlalu rendah.  Buatlah sendiri perkiraan yang masuk akal untuk distribusi $m(x)$, lalu tentukan
peluang bahwa rata-rata nilai seorang mahasiswa akurat hingga .05.  Tentukan juga peluang bahwa
rata-rata tersebut akurat hingga .5.

\i\label{exer 9.2.114}  Carilah barisan peubah acak diskret independen terbatas seragam $\{X_n\}$
sedemikian sehingga varians jumlahnya tidak menuju $\infty$ ketika $n \rightarrow \infty$, dan
jumlahnya tidak berdistribusi normal secara asimtotik.

\end{LJSItem}

% As a precaution, I commented this choice command out.
%\choice{}{\section[Continuous Independent Trials]{Central Limit Theorem for 
\section[Percobaan Independen Kontinu]{Teorema Limit Pusat untuk
Percobaan Independen Kontinu}
\label{sec 9.4}
Dalam Bagian~\ref{sec 9.3}, kita telah melihat bahwa fungsi distribusi bagi jumlah dari banyak
peubah acak diskret independen, yang banyaknya dinyatakan oleh $n$, dengan rataan~$\mu$ dan
varians~$\sigma^2$ cenderung menyerupai densitas normal dengan rataan~$n\mu$ dan
varians~$n\sigma^2$.  Hal yang luar biasa ialah bahwa hasil ini berlaku bagi {\em setiap}
distribusi dengan rataan dan varians berhingga.  Dalam bagian ini, kita akan melihat bahwa hasil
yang sama juga berlaku bagi peubah acak
kontinu yang memiliki fungsi densitas yang sama.
\par
Mari kita mulai dengan mengamati beberapa contoh untuk melihat apakah hasil semacam itu masuk akal.

\subsection*{Jumlah Terstandardisasi}
\begin{example}
Andaikan kita memilih $n$ bilangan acak dari interval $[0,1]$ dengan densitas seragam.  Misalkan
$X_1$,~$X_2$, \dots,~$X_n$ menyatakan pilihan-pilihan tersebut, dan $S_n = X_1 + X_2 +\cdots+ X_n$
menyatakan jumlahnya.

Dalam Contoh~\ref{exam 7.12}, kita melihat bahwa fungsi densitas~$S_n$ cenderung berbentuk normal,
tetapi berpusat di~$n/2$ dan semakin mendatar.  Untuk membandingkan bentuk fungsi-fungsi densitas
ini bagi nilai~$n$ yang berbeda, kita mengikuti cara dalam seksi sebelumnya: kita
\emx {menstandardisasi} $S_n$ dengan mendefinisikan
$$
S_n^* = \frac {S_n - n\mu}{\sqrt n \sigma}\ .
$$
Dengan demikian, untuk semua~$n$ kita memperoleh
\begin{eqnarray*}
E(S_n^*) & = & 0\ , \\
V(S_n^*) & = & 1\ .
\end{eqnarray*}
Fungsi densitas~$S_n^*$ hanyalah versi terstandardisasi dari fungsi densitas~$S_n$ (lihat
Gambar~\ref{fig 9.7}).
\end{example}

\putfig{4truein}{PSfig9-7}
{Fungsi densitas $S^*_n$ (kasus seragam, $n = 2, 3, 4, 10$).}{fig 9.7} 

\begin{example}
Mari kita melakukan hal yang sama, tetapi sekarang kita memilih bilangan dari interval
$[0,+\infty)$ dengan densitas eksponensial berparameter~$\lambda$.  Maka (lihat
Contoh~\ref{exam 6.21})
\par
\begin{eqnarray*}
\mu & = & E(X_i)  =  \frac 1\lambda\ , \\
\sigma^2 & = & V(X_j) = \frac 1{\lambda^2}\ .
\end{eqnarray*}
\par
Di sini kita mengetahui fungsi densitas~$S_n$ secara eksplisit (lihat
Bagian~\ref{sec 7.2}).  Kita dapat menggunakan Akibat~\ref{cor 5.1} untuk menghitung fungsi
densitas $S_n^*$.  Kita memperoleh
\par
\begin{eqnarray*}
  f_{S_n}(x) & = & \frac {\lambda e^{-\lambda x}(\lambda x)^{n - 1}}{(n - 1)!}\ , \\
f_{S_n^*}(x) & = & \frac {\sqrt n}\lambda f_{S_n} \left( \frac {\sqrt n x +
n}\lambda \right)\ .
\end{eqnarray*}
Grafik fungsi densitas~$S_n^*$ ditampilkan dalam Gambar~\ref{fig 9.9}.
\end{example}

\putfig{4truein}{PSfig9-9}
{Fungsi densitas $S^*_n$ (kasus eksponensial,
$n = 2, 3, 10, 30$, $\lambda = 1$).}{fig 9.9} 

Contoh-contoh ini membuat masuk akal dugaan bahwa fungsi densitas peubah acak yang dinormalisasi
$S_n^*$, untuk~$n$ yang besar, akan sangat menyerupai densitas normal dengan rataan~0 dan varians~1, baik
dalam kasus kontinu maupun kasus diskret.  Teorema Limit Pusat menyatakan dugaan ini secara tepat.

\subsection*{Teorema Limit Pusat}
\begin{theorem}\label{thm 9.4.7}{\bf (Teorema Limit Pusat)}\index{Teorema Limit Pusat!untuk
proses percobaan independen kontinu} 
Misalkan $S_n = X_1 + X_2 +\cdots+ X_n$ adalah jumlah $n$ peubah acak kontinu independen dengan
fungsi densitas yang sama~$p$, nilai harapan~$\mu$, dan varians~$\sigma^2$.  Misalkan
$S_n^* = (S_n - n\mu)/\sqrt n \sigma$.  Maka, untuk semua $a < b$, kita mempunyai
$$
\lim_{n \to \infty} P(a < S_n^* < b) = \frac 1{\sqrt{2\pi}} \int_a^b
e^{-x^2/2}\, dx\ .
$$
\end{theorem}
%********Some applications of the CLT to statistics should be given 
%at the end of Section 9.2.  We should also change the structure of
%the examples here relative to the subsections. 
\par
Kita akan memberikan bukti teorema ini dalam Bagian~\ref{sec 10.3}.  Sekarang kita akan membahas
beberapa contoh.

\begin{example}\label{exam 9.10}
Andaikan seorang juru ukur ingin mengukur jarak yang telah diketahui, misalnya 1~mil, dengan
menggunakan alat transit dan suatu metode triangulasi.  Ia mengetahui bahwa akibat kemungkinan
pergeseran alat transit, distorsi atmosfer, dan kesalahan manusia, setiap pengukuran cenderung
sedikit meleset.  Ia berencana melakukan beberapa pengukuran dan mengambil rata-ratanya.  Ia
mengandaikan bahwa pengukuran-pengukurannya merupakan peubah acak independen dengan distribusi yang
sama, rataan $\mu = 1$, dan simpangan baku $\sigma = .0002$ (jadi, jika galat berdistribusi
mendekati normal, pengukurannya berada dalam jarak 1 kaki dari jarak yang benar sekitar 65\% dari
seluruh percobaan).  Apa yang dapat ia katakan mengenai rata-ratanya?
\par
Ia dapat mengatakan bahwa jika $n$ besar, rata-rata $S_n/n$ memiliki fungsi densitas yang mendekati
normal, dengan rataan $\mu = 1$ mil dan simpangan baku $\sigma = .0002/\sqrt n$ mil.

Berapa banyak pengukuran yang perlu ia lakukan agar cukup yakin bahwa rata-ratanya berada dalam
jarak .0001 dari nilai sebenarnya?  Ketaksamaan Chebyshev menyatakan
$$
P\left(\left| \frac {S_n}n - \mu \right| \geq .0001 \right) \leq \frac
{(.0002)^2}{n(10^{-8})} = \frac 4n\ ,
$$
sehingga kita harus memiliki $n \ge 80$ agar peluang bahwa galatnya kurang dari .0001 melampaui
.95.
\par
Kita telah memperhatikan bahwa perkiraan dari ketaksamaan Chebyshev tidak selalu baik, dan inilah
salah satu contohnya.  Jika kita mengandaikan bahwa $n$ cukup besar sehingga densitas~$S_n$
mendekati normal, maka kita mempunyai
\par
\begin{eqnarray*}
P\left(\left| \frac {S_n}n - \mu \right| < .0001 \right) &=& P\bigl(-.5\sqrt{n} < S_n^*
< +.5\sqrt{n}\bigr) \\
     &\approx& \frac 1{\sqrt{2\pi}} \int_{-.5\sqrt{n}}^{+.5\sqrt{n}} e^{-x^2/2}\, dx\ ,
\end{eqnarray*}
dan pernyataan terakhir ini lebih besar dari .95 jika $.5\sqrt{n} \ge 2.$  Artinya, cukup mengambil
$n = 16$ pengukuran untuk memperoleh hasil yang sama.  Perhitungan kedua ini lebih tajam, tetapi
bergantung pada asumsi bahwa $n = 16$ cukup besar agar densitas normal menjadi aproksimasi yang baik
bagi~$S_n^*$, dan dengan demikian bagi~$S_n$.  Di sini Teorema Limit Pusat tidak mengatakan seberapa
besar $n$ harus dipilih.  Dalam kebanyakan kasus yang melibatkan jumlah peubah acak independen,
pedoman praktis yang baik ialah bahwa aproksimasi sudah baik untuk $n \ge 30$.  Dalam kasus ini,
jika kita mengandaikan bahwa galat berdistribusi mendekati normal, aproksimasinya mungkin cukup baik
bahkan untuk $n = 16$.
\end{example}

\subsection*{Memperkirakan Rataan}

\begin{example}(Lanjutan Contoh~\ref{exam 9.10})
Sekarang andaikan juru ukur kita mengukur jarak yang belum diketahui dengan peralatan yang sama
dalam kondisi yang sama.  Ia melakukan 36 pengukuran dan menghitung rata-ratanya.  Seberapa yakin
ia dapat menyatakan bahwa hasil pengukurannya berada dalam jarak .0002 dari nilai sebenarnya?

Dengan kembali menggunakan aproksimasi normal, kita memperoleh
\begin{eqnarray*}
P\left(\left|\frac {S_n}n - \mu\right| < .0002 \right) &=& P\bigl(|S_n^*| < .5\sqrt n\bigr) \\
     &\approx& \frac 2{\sqrt{2\pi}} \int_{-3}^3 e^{-x^2/2}\, dx \\
     &\approx& .997\ .
\end{eqnarray*}


Ini berarti bahwa juru ukur tersebut dapat memiliki keyakinan sebesar 99.7~persen bahwa rata-ratanya berada
dalam jarak .0002 dari nilai sebenarnya.  Untuk meningkatkan tingkat keyakinannya, ia dapat
melakukan lebih banyak pengukuran, menuntut ketelitian yang lebih rendah, atau meningkatkan mutu
pengukurannya (yaitu mengurangi varians~$\sigma^2$).  Dalam setiap kasus, Teorema Limit Pusat
memberikan informasi kuantitatif mengenai tingkat keyakinan suatu proses pengukuran, dengan tetap
mengandaikan bahwa aproksimasi normal berlaku.
\par
Sekarang andaikan juru ukur tersebut tidak mengetahui rataan ataupun simpangan baku
pengukuran-pengukurannya, tetapi mengandaikan bahwa pengukuran itu independen.  Apa yang harus ia
lakukan?
\par
Sekali lagi, ia melakukan beberapa pengukuran terhadap jarak yang telah diketahui dan menghitung
rata-ratanya.  Seperti sebelumnya, rata-rata galat berdistribusi mendekati normal, tetapi sekarang
dengan rataan dan varians yang tidak diketahui.
\end{example}
\subsection*{Rataan Sampel}
Jika ia mengetahui bahwa varians~$\sigma^2$ distribusi galat adalah .0002, ia dapat memperkirakan
rataan~$\mu$ dengan mengambil \emx {rata-rata,} atau \emx {rataan sampel} dari, misalnya, 36
pengukuran:
$$
\bar \mu = \frac {x_1 + x_2 +\cdots+ x_n}n\ ,
$$
dengan $n = 36$.  Kemudian, seperti sebelumnya, $E(\bar \mu) = \mu$.  Selain itu, argumen di atas
menunjukkan bahwa
$$
P(|\bar \mu - \mu| < .0002) \approx .997\ .
$$
Interval $(\bar \mu - .0002, \bar \mu + .0002)$ disebut \emx {interval kepercayaan
99.7\%}\index{interval kepercayaan} untuk~$\mu$ (lihat Contoh~\ref{exam 9.4.1}).

\subsection*{Varians Sampel}
Jika ia tidak mengetahui varians~$\sigma^2$ distribusi galat, ia dapat memperkirakan $\sigma^2$
dengan \emx {varians sampel}:
$$
\bar \sigma^2 = \frac {(x_1 - \bar \mu)^2 + (x_2 - \bar \mu)^2
+\cdots+ (x_n - \bar \mu)^2}n\ ,
$$
dengan $n = 36$.  Hukum Bilangan Besar, ketika diterapkan pada peubah-peubah acak
$(X_i - \bar \mu)^2$, menyatakan bahwa untuk~$n$ yang besar, varians sampel~$\bar \sigma^2$
mendekati varians~$\sigma^2$, sehingga juru ukur dapat menggunakan $\bar \sigma^2$ sebagai
pengganti~$\sigma^2$ dalam argumen di atas.
\par
Pengalaman menunjukkan bahwa, dalam sebagian besar masalah praktis jenis ini, varians sampel
merupakan perkiraan yang baik bagi varians dan dapat digunakan sebagai pengganti varians untuk
menentukan tingkat kepercayaan rataan sampel.  Artinya, kita dapat mengandalkan Hukum Bilangan Besar
untuk memperkirakan varians dan Teorema Limit Pusat untuk memperkirakan rataan.
\par
Kita dapat memeriksa hal ini dalam beberapa kasus khusus.  Andaikan kita mengetahui bahwa distribusi
galat adalah \emx {normal,} dengan rataan dan varians yang tidak diketahui.  Kita dapat mengambil
sampel yang terdiri dari $n$ pengukuran, mencari rataan sampel~$\bar \mu$ dan varians
sampel~$\bar \sigma^2$, lalu membentuk
$$
T_n^* = \frac {S_n - n\bar\mu}{\sqrt{n}\bar\sigma}\ ,
$$
dengan $n = 36$.  Kita mengharapkan $T_n^*$ menjadi aproksimasi yang baik bagi~$S_n^*$ untuk~$n$
yang besar.

\subsection*{Densitas-$t$}
Statistikawan W.~S. Gosset\index{GOSSET, W. S.}\footnote{W. S. Gosset menemukan distribusi yang
sekarang kita sebut distribusi-$t$ ketika bekerja untuk Guinness Brewery di Dublin.  Ia menulis
dengan nama samaran ``Student."  Hasil yang dibahas di sini pertama kali terbit dalam Student,
``The Probable Error of a Mean," \emx {Biometrika,} vol.~6 (1908), pp.~1-24.} menunjukkan bahwa
dalam kasus ini $T_n^*$ memiliki fungsi densitas yang bukan normal, melainkan
\emx {densitas-$t$}\index{fungsi densitas!t-}\index{densitas-t} dengan $n$ derajat kebebasan.
(Jumlah derajat kebebasan, yaitu $n$, hanyalah parameter yang memberi tahu kita densitas-$t$ mana yang
harus digunakan.)  Dalam kasus ini, kita dapat menggunakan densitas-$t$ sebagai pengganti densitas
normal untuk menentukan tingkat kepercayaan bagi~$\mu$.  Ketika $n$ bertambah, densitas-$t$
mendekati densitas normal.  Bahkan, untuk $n = 8$, densitas-$t$ dan densitas normal praktis sama
(lihat Gambar~\ref{fig 9.12}).

\putfig{4.5truein}{PSfig9-12}
{Grafik $t-$densitas untuk $n= 1, 3, 8$ dan densitas normal dengan $\mu = 0,
\sigma = 1$.}{fig 9.12} 

\exercises
\indent \emx {Catatan untuk soal komputer}:
\begin{description}
\item[(a)] $\ $Simulasi: Ingat (lihat Akibat~\ref{cor 5.2}) bahwa
$$
X = F^{-1}(rnd)   
$$
akan menyimulasikan peubah acak dengan densitas $f(x)$ dan distribusi
$$
F(X) = \int_{-\infty}^x f(t)\, dt\ .
$$
Jika $f(x)$ merupakan fungsi densitas normal dengan rataan $\mu$ dan
simpangan baku $\sigma$, sedangkan baik
$F$ maupun $F^{-1}$ tidak dapat
dinyatakan dalam bentuk tertutup, gunakan sebagai gantinya
$$
X = \sigma\sqrt {-2\log(rnd)} \cos 2\pi(rnd) + \mu\ .
$$
\item[(b)] $\ $Grafik batang: upayakan agar grafik Anda memuat sekitar 20~hingga~30 batang
(dengan lebar yang sama).  Hal ini dapat dicapai dengan memilih rentang $[x{\rm min}, x{\rm min}]$
dan banyaknya batang secara tepat (misalnya, $[\mu - 3\sigma, \mu + 3\sigma]$ dengan 30 batang
akan cocok dalam banyak kasus).  Bereksperimenlah!
\end{description}
\vskip .1in

\begin{LJSItem}

\i\label{exer 9.4.1}  Misalkan $X$ adalah peubah acak kontinu dengan rataan $\mu(X)$ dan varians
$\sigma^2(X)$, dan misalkan $X^* = (X - \mu)/\sigma$ adalah versi terstandardisasinya.
Verifikasikan secara langsung bahwa $\mu(X^*) = 0$ dan $\sigma^2(X^*) = 1$.

\i\label{exer 9.4.2}  Misalkan $\{X_k\}$, $1 \leq k \leq n$, adalah barisan peubah acak
independen, semuanya dengan rataan~0 dan varians~1, dan misalkan $S_n$,~$S_n^*$, dan~$A_n$
berturut-turut adalah jumlah, jumlah terstandardisasi, dan rata-ratanya.  Verifikasikan secara
langsung bahwa
$S_n^* = S_n/\sqrt{n} = \sqrt{n} A_n$.

\i\label{exer 9.4.3}  Misalkan $\{X_k\}$, $1 \leq k \leq n$, adalah barisan peubah acak yang
semuanya memiliki rataan~$\mu$ dan varians~$\sigma^2$, dan misalkan $Y_k = X_k^*$ adalah versi
terstandardisasi dari peubah-peubah tersebut.  Misalkan $S_n$ dan $T_n$ masing-masing adalah jumlah
peubah-peubah $X_k$ dan $Y_k$, sedangkan $S_n^*$ dan $T_n^*$ adalah versi terstandardisasi dari
jumlah-jumlah itu.  Tunjukkan bahwa $S_n^* = T_n^* =
T_n/\sqrt{n}$.

\i\label{exer 9.4.4} Andaikan kita memilih secara independen 25 bilangan acak
(dengan densitas seragam) dari interval $[0,20]$.  Tuliskan densitas-densitas normal yang
mengaproksimasi densitas jumlahnya $S_{25}$, jumlah terstandardisasinya
$S_{25}^*$, dan rata-ratanya $A_{25}$.

\i\label{exer 9.4.5} Tulislah program untuk memilih secara independen 25 bilangan acak dari
$[0,20]$, menghitung jumlahnya $S_{25}$, dan mengulangi percobaan ini 1000 kali.  Buatlah grafik
batang untuk densitas~$S_{25}$ dan bandingkan dengan aproksimasi normal dalam
Latihan~\ref{exer 9.4.4}.  Seberapa baik kecocokannya?  Sekarang lakukan hal yang sama untuk jumlah
terstandardisasi $S_{25}^*$ dan rata-rata $A_{25}$.

\i\label{exer 9.4.6}  Secara umum, Teorema Limit Pusat memberikan perkiraan yang lebih baik daripada
ketaksamaan Chebyshev untuk rata-rata suatu jumlah.  Untuk melihatnya, misalkan $A_{25}$ adalah
rata-rata yang dihitung dalam Latihan~\ref{exer 9.4.5}, dan misalkan $N$ adalah aproksimasi normal
bagi~$A_{25}$.  Ubah program Anda dalam Latihan~\ref{exer 9.4.5} agar menghasilkan tabel fungsi
$F(x) = P(|A_{25} - 10| \geq x) = {}$ proporsi dari keseluruhan 1000 percobaan yang memenuhi
$|A_{25} - 10| \geq x$.  Lakukan hal yang sama untuk fungsi $f(x) = P(|N - 10| \geq x)$.
(Untuk ini, Anda dapat menggunakan tabel normal, Tabel~\ref{tabl 9.1}, atau prosedur
{\bf NormalArea}.)  Sekarang gambarkan pada sumbu yang sama grafik $F(x)$,~$f(x)$, dan fungsi
Chebyshev $g(x) = 4/(3x^2)$.  Bagaimana perbandingan $f(x)$ dan $g(x)$ sebagai perkiraan bagi
$F(x)$?

\i\label{exer 9.4.7} Teorema Limit Pusat menyatakan bahwa jumlah peubah-peubah acak independen
cenderung tampak normal, betapa pun tidak lazimnya distribusi setiap peubah itu.  Mari kita mengujinya
dengan simulasi komputer.  Pilihlah secara independen 25 bilangan dari interval $[0,1]$ dengan
densitas peluang $f(x)$ yang diberikan di bawah, lalu hitung jumlahnya $S_{25}$.  Ulangi percobaan ini
1000 kali dan buatlah grafik batang dari hasilnya.  Sekarang gambarkan pada grafik yang sama densitas
$\phi(x) = \mbox {normal \,\,\,}(x,\mu(S_{25}),\sigma(S_{25}))$.  Seberapa baik densitas normal
cocok dengan grafik batang Anda dalam setiap kasus?
\begin{enumerate}
\item  $f(x) = 1$.

\item  $f(x) = 2x$.

\item  $f(x) = 3x^2$.

\item  $f(x) = 4|x - 1/2|$.

\item  $f(x) = 2 - 4|x - 1/2|$.
\end{enumerate}

\i\label{exer 9.4.8}  Ulangi percobaan yang dijelaskan dalam Latihan~\ref{exer 9.4.7}, tetapi
sekarang pilihlah 25 bilangan dari $[0,\infty)$ dengan menggunakan $f(x) = e^{-x}$.

\i\label{exer 9.4.9}  Seberapa besar $n$ harus dipilih sebelum
$S_n = X_1 + X_2 +\cdots+ X_n$ mendekati distribusi normal?  Bilangan ini sering kali ternyata
sangat kecil.  Mari kita selidiki pertanyaan ini dengan simulasi komputer.  Pilihlah $n$ bilangan
dari $[0,1]$ dengan densitas peluang $f(x)$, dengan $n = 3$,~6, 12,~20, dan ambil $f(x)$ secara
bergiliran sebagai setiap densitas dalam Latihan~\ref{exer 9.4.7}.  Hitung jumlahnya $S_n$, ulangi
percobaan ini 1000 kali, dan buatlah grafik batang hasilnya yang terdiri dari 20 batang.  Seberapa
besar $n$ harus dipilih sebelum diperoleh kecocokan yang baik?

\i\label{exer 9.4.10}  Seorang juru ukur sedang mengukur tinggi tebing yang diketahui sekitar
1000 kaki.  Ia mengandaikan bahwa alatnya telah dikalibrasi dengan benar dan galat-galat pengukurannya
independen, dengan rataan $\mu = 0$ dan varians $\sigma^2 = 10$.  Ia berencana melakukan $n$
pengukuran lalu menghitung rata-ratanya.  Dengan menggunakan (a)~ketaksamaan Chebyshev dan
(b)~aproksimasi normal, perkirakan seberapa besar $n$ harus dipilih jika ia ingin 95~persen yakin
bahwa rata-ratanya berada dalam jarak 1~kaki dari nilai sebenarnya.  Sekarang, dengan menggunakan
(a) dan (b), perkirakan berapa nilai $\sigma^2$ yang diperlukan jika ia hanya ingin melakukan
10 pengukuran dengan tingkat kepercayaan yang sama.

\i\label{exer 9.4.11} Harga satu lembar saham Pilsdorff Beer Company (lihat
Latihan~\ref{sec 8.2}.\ref{exer 8.2.12}) dinyatakan oleh $Y_n$ pada hari ke-$n$ dalam setahun.
Finn mengamati bahwa selisih $X_n = Y_{n + 1} - Y_n$ tampaknya merupakan peubah acak independen
dengan distribusi yang sama, yang memiliki rataan $\mu = 0$ dan varians $\sigma^2 = 1/4$.
Jika $Y_1 = 100$, perkirakan peluang bahwa $Y_{365}$
\begin{enumerate}
\item  ${} \geq 100$.

\item  ${} \geq 110$.

\item  ${} \geq 120$.
\end{enumerate}

\i\label{exer 9.4.12}  Ujilah kesimpulan Anda dalam Latihan~\ref{exer 9.4.11} dengan simulasi
komputer.  Pertama-tama, pilihlah 364 bilangan $X_i$ dengan densitas
$f(x) = \mbox {normal}(x,0,1/4)$.  Kemudian bentuk jumlah
$Y_{365} = 100 + X_1 + X_2 +\cdots+ X_{364}$ dan ulangi percobaan ini 200 kali.  Buatlah grafik
batang hasilnya pada $[50,150]$, dengan menumpangkan grafik densitas normal yang menjadi
aproksimasinya.  Apa yang ditunjukkan grafik ini mengenai jawaban Anda dalam
Latihan~\ref{exer 9.4.11}?

\i\label{exer 9.4.12.5}  Fisikawan menyatakan bahwa partikel-partikel dalam tabung panjang terus
bergerak bolak-balik sepanjang tabung, masing-masing dengan kecepatan $V_k$ (dalam cm/detik) yang
pada setiap saat tertentu berdistribusi normal, dengan rataan $\mu = 0$ dan varians
$\sigma^2 = 1$.  Andaikan terdapat $10^{20}$ partikel dalam tabung tersebut.
\begin{enumerate}
\item  Carilah rataan dan varians kecepatan rata-rata partikel-partikel itu.

\item  Berapa peluang bahwa kecepatan rata-ratanya ${} \geq
10^{-9}$~cm/detik?
\end{enumerate}

\i\label{exer 9.4.13}  Seorang astronom melakukan $n$ pengukuran jarak antara Jupiter dan salah
satu satelitnya.  Pengalaman menggunakan alat-alat tersebut membuatnya yakin bahwa, dengan satuan
yang sesuai, pengukuran-pengukuran itu akan berdistribusi normal dengan rataan~$d$, yaitu jarak
sebenarnya, dan varians~16.  Ia melakukan serangkaian $n$ pengukuran.  Misalkan
$$
A_n = \frac {X_1 + X_2 +\cdots+ X_n}n
$$
adalah rata-rata pengukuran-pengukuran tersebut.
\begin{enumerate}
\item  Tunjukkan bahwa
\[
P\left(A_n - \frac 8{\sqrt n} \leq d \leq A_n + \frac 8{\sqrt n}\right) \approx
.95.
\]

\item  Ketika sembilan pengukuran dilakukan, rata-rata jaraknya ternyata 23.2~satuan.
Memasukkan nilai yang diamati ke dalam (a) memberikan \emx {interval kepercayaan 95~persen}
bagi jarak~$d$ yang belum diketahui.  Hitunglah interval ini.

\item  Mengapa dalam (b) kita tidak menyatakan secara lebih sederhana bahwa peluang nilai~$d$
terletak dalam interval kepercayaan yang dihitung adalah .95?

\item  Perubahan apa yang akan Anda buat dalam prosedur di atas jika Anda ingin menghitung
interval kepercayaan 99~persen?
\end{enumerate}

\i\label{exer 9.4.14}  Gambarkan grafik batang yang serupa dengan grafik dalam
Gambar~\ref{fig 9.61} untuk tinggi rerata orang tua dalam data Galton sebagaimana diberikan dalam
Lampiran~B, lalu bandingkan grafik batang ini dengan kurva normal yang sesuai.

\end{LJSItem}
%\end{LJSItem}}
%\end{document}
